% 机器人 VLA 实时运行时：汇报稿（16:9 幻灯片）
% 对象：Pi0.5 的实时输入与推理、三框架的能力、Level 1 测量、Level 2 调研与待决定事项。
% 编译：bash build.sh。共用的字体、配色、TikZ 样式与宏在 common.tex，幻灯片版式在 deck.tex。
% 数字来源：docs/design.md 与 runs/ 下的结果文件；每页的来源在页脚注记中给出。
\PassOptionsToPackage{table}{xcolor}
\documentclass[aspectratio=169,10pt]{beamer}

% 两份文档共用的部分：字体、配色、TikZ 样式、表格与代码块、标签与内容宏。
% 不含页面版式：幻灯片外观在 deck.tex，两份文档的标题与内容在 main.tex、main_framework.tex。
% 使用前需要先在调用文件里写 \PassOptionsToPackage{table}{xcolor} 和 \documentclass。

% ---------- 字体：Lato（拉丁）+ Noto Sans CJK SC（中文），等宽用 DejaVu Sans Mono ----------
\usepackage[UTF8,scheme=plain]{ctex}
% 注意：fontconfig 会把名称 "Lato" 解析成 Hairline 字重，所以直接按文件加载；缺文件时退回 Noto Sans CJK SC。
\IfFileExists{/usr/share/fonts/truetype/lato/Lato-Regular.ttf}{%
  \setmainfont{Lato-Regular.ttf}[Path=/usr/share/fonts/truetype/lato/,BoldFont=Lato-Bold.ttf,ItalicFont=Lato-Italic.ttf,BoldItalicFont=Lato-BoldItalic.ttf]%
  \setsansfont{Lato-Regular.ttf}[Path=/usr/share/fonts/truetype/lato/,BoldFont=Lato-Bold.ttf,ItalicFont=Lato-Italic.ttf,BoldItalicFont=Lato-BoldItalic.ttf]%
}{%
  \setmainfont{Noto Sans CJK SC}\setsansfont{Noto Sans CJK SC}%
}
\setmonofont{DejaVu Sans Mono}[Scale=0.84]
\setCJKmainfont{Noto Sans CJK SC}
\setCJKsansfont{Noto Sans CJK SC}
\setCJKmonofont{Noto Sans CJK SC}
\renewcommand{\familydefault}{\sfdefault}

\usepackage{booktabs,array,tabularx,multirow,graphicx,xcolor}
\usepackage{tikz}
\usetikzlibrary{arrows.meta,positioning,calc,fit,shapes.geometric,backgrounds,decorations.pathreplacing}
\usepackage{fancyvrb}
\usepackage[skins]{tcolorbox}
\usepackage{hyperref}

% ---------- 配色：取自公司模板（深蓝 001E50、蓝 0054AE / 0046C9 / 00377C、青 00C7FD、紫 6515AD / AF99FF、浅底 F5F8FF）----------
% 语义：Teal = M*（主蓝）、Orange = vLLM-Omni（紫）、Blue = SGLang-Omni（青）、Gold = 状态 / 交接（深蓝）、Red = 警示。
% 变量名沿用旧版（Teal / Orange / ...），只改取值，页面内容里的颜色引用不用动。
\definecolor{Ink}{HTML}{0A1A33}
\definecolor{Navy}{HTML}{001E50}
\definecolor{Teal}{HTML}{0054AE}
\definecolor{Blue}{HTML}{0092C2}
\definecolor{Orange}{HTML}{6515AD}
\definecolor{Red}{HTML}{C2334A}
\definecolor{Gold}{HTML}{00377C}
\definecolor{Cyan}{HTML}{00C7FD}
\definecolor{Lav}{HTML}{AF99FF}
\definecolor{Muted}{HTML}{5B6B82}
\definecolor{Line}{HTML}{C9D3E3}
\definecolor{Mist}{HTML}{F5F8FF}
\definecolor{TealBg}{HTML}{E6EFF9}
\definecolor{BlueBg}{HTML}{E0F6FD}
\definecolor{OrangeBg}{HTML}{EFE7F8}
\definecolor{RedBg}{HTML}{FBE9EC}
\definecolor{GoldBg}{HTML}{E4EAF4}
\hypersetup{colorlinks=true,linkcolor=Ink,urlcolor=Blue,citecolor=Blue}

\newcolumntype{Y}{>{\raggedright\arraybackslash}X}
% Pull top-aligned columns closer to the title rule without changing table-only slides.
\newcommand{\minipagetop}{\vspace{-.3cm}}

% ---------- 通用宏 ----------
\newcommand{\memo}[1]{{\color{Muted}\footnotesize #1}}
\newcommand{\pill}[2]{{\setlength{\fboxsep}{1.6pt}\raisebox{.05ex}{\colorbox{#1}{\textcolor{white}{\fontsize{6.5}{7}\selectfont\bfseries #2}}}}}
\newcommand{\tagE}{\pill{Teal}{实验}}
\newcommand{\tagC}{\pill{Blue}{源码}}
\newcommand{\tagO}{\pill{Gold}{官方}}
\newcommand{\tagH}{\pill{Orange}{假设}}
\newcommand{\tagG}{\pill{Red}{未测}}
\newcommand{\labSame}{\pill{Teal}{方向一致}}
\newcommand{\labMag}{\pill{Gold}{量级不同}}
\newcommand{\labBase}{\pill{Orange}{口径不同}}
\newcommand{\labRev}{\pill{Red}{方向相反}}
% 伪代码块：Verbatim 中用 \kw{} \cm{} \id{} 着色（花括号需写成 \{ \}）
\newcommand{\kw}[1]{\textbf{\textcolor{Teal}{#1}}}
\newcommand{\cm}[1]{\textcolor{Muted}{#1}}
\newcommand{\id}[1]{\textcolor{Orange}{#1}}
\newcommand{\hi}[1]{\textbf{\textcolor{Red}{#1}}}
\newtcolorbox{codebox}[2][]{enhanced,colback=Mist,colframe=Line,boxrule=.4pt,arc=2pt,left=4pt,right=3pt,top=2pt,bottom=2pt,
  title={#2},coltitle=white,colbacktitle=Teal,fonttitle=\fontsize{6.8}{8}\selectfont\bfseries,boxsep=1pt,toptitle=1.2pt,bottomtitle=1.2pt,lefttitle=4pt,before skip=0pt,#1}
% 右侧说明框：高度自适应，可连续堆叠
\newtcolorbox{nbox}[3][]{colback=#2!6,colframe=#2,boxrule=.7pt,arc=2.5pt,left=4pt,right=4pt,top=3pt,bottom=3pt,
  before upper={\def\codesize{\relax}\noindent\textbf{\textcolor{#2}{#3}}\par\vspace{1pt}},fontupper=\fontsize{7.7}{9.5}\selectfont\raggedright,before skip=0pt,after skip=4pt,#1}
\newcommand{\hl}[1]{\textbf{\textcolor{Teal}{#1}}}
\newcommand{\warn}[1]{\textbf{\textcolor{Red}{#1}}}
\newcommand{\codesize}{\small}
\newcommand{\code}[1]{\texttt{\codesize #1}}
\newcommand{\tb}[1]{{\bfseries #1}}
% Walk 名字：小字号，长名（如 image_gen_guided3）能放进 2.25 cm 的框
\newcommand{\wname}[1]{{\fontsize{6.6}{8}\selectfont\bfseries #1}}

% 图片来源注记
\newcommand{\figsource}[1]{{\color{Muted}\fontsize{6.5}{7.8}\selectfont #1}}

% ---------- TikZ 图元样式（固定高度、字号，避免重叠） ----------
\tikzset{
  >={Latex[length=1.8mm]},
  box/.style={draw=#1,fill=#1!8,rounded corners=2.5pt,line width=.8pt,align=center,font=\footnotesize,minimum height=.62cm,inner xsep=4pt,inner ysep=2.5pt},
  box/.default=Teal,
  sbox/.style={box=#1,font=\scriptsize,minimum height=.45cm,inner xsep=3pt,inner ysep=1.5pt},
  sbox/.default=Teal,
  plain/.style={align=center,font=\footnotesize,inner sep=1pt},
  tiny/.style={align=center,font=\scriptsize,inner sep=1pt,text=Muted},
  flow/.style={->,draw=Muted,line width=.9pt},
  ctrl/.style={->,draw=Blue,line width=.9pt,densely dashed},
  strm/.style={->,draw=Orange,line width=1.2pt},
  back/.style={->,draw=Teal,line width=1pt,rounded corners=3pt},
  wk/.style={box=#1,minimum height=.62cm,font=\scriptsize,inner xsep=3pt},
  panel/.style={draw=Line,fill=Mist,rounded corners=3pt,line width=.6pt},
}

% 对数轴上的点与连线（第 18 页）
\newcommand{\two}[3]{\pgfmathsetmacro{\xa}{ln(#2/0.5)/ln(2)*1.5}\pgfmathsetmacro{\xb}{ln(#3/0.5)/ln(2)*1.5}\draw[Muted!60,line width=.8pt] (\xa,#1)--(\xb,#1);}
\newcommand{\dotat}[3]{\pgfmathsetmacro{\xx}{ln(#2/0.5)/ln(2)*1.5}\node[#3] at (\xx,#1) {};}

% 幻灯片版式（两份汇报共用）：16:9 Beamer，白底，左对齐标题，底部固定“要点条”，右下页码。
% 调用前需要：\documentclass[aspectratio=169,10pt]{beamer}，然后 % 两份文档共用的部分：字体、配色、TikZ 样式、表格与代码块、标签与内容宏。
% 不含页面版式：幻灯片外观在 deck.tex，两份文档的标题与内容在 main.tex、main_framework.tex。
% 使用前需要先在调用文件里写 \PassOptionsToPackage{table}{xcolor} 和 \documentclass。

% ---------- 字体：Lato（拉丁）+ Noto Sans CJK SC（中文），等宽用 DejaVu Sans Mono ----------
\usepackage[UTF8,scheme=plain]{ctex}
% 注意：fontconfig 会把名称 "Lato" 解析成 Hairline 字重，所以直接按文件加载；缺文件时退回 Noto Sans CJK SC。
\IfFileExists{/usr/share/fonts/truetype/lato/Lato-Regular.ttf}{%
  \setmainfont{Lato-Regular.ttf}[Path=/usr/share/fonts/truetype/lato/,BoldFont=Lato-Bold.ttf,ItalicFont=Lato-Italic.ttf,BoldItalicFont=Lato-BoldItalic.ttf]%
  \setsansfont{Lato-Regular.ttf}[Path=/usr/share/fonts/truetype/lato/,BoldFont=Lato-Bold.ttf,ItalicFont=Lato-Italic.ttf,BoldItalicFont=Lato-BoldItalic.ttf]%
}{%
  \setmainfont{Noto Sans CJK SC}\setsansfont{Noto Sans CJK SC}%
}
\setmonofont{DejaVu Sans Mono}[Scale=0.84]
\setCJKmainfont{Noto Sans CJK SC}
\setCJKsansfont{Noto Sans CJK SC}
\setCJKmonofont{Noto Sans CJK SC}
\renewcommand{\familydefault}{\sfdefault}

\usepackage{booktabs,array,tabularx,multirow,graphicx,xcolor}
\usepackage{tikz}
\usetikzlibrary{arrows.meta,positioning,calc,fit,shapes.geometric,backgrounds,decorations.pathreplacing}
\usepackage{fancyvrb}
\usepackage[skins]{tcolorbox}
\usepackage{hyperref}

% ---------- 配色：取自公司模板（深蓝 001E50、蓝 0054AE / 0046C9 / 00377C、青 00C7FD、紫 6515AD / AF99FF、浅底 F5F8FF）----------
% 语义：Teal = M*（主蓝）、Orange = vLLM-Omni（紫）、Blue = SGLang-Omni（青）、Gold = 状态 / 交接（深蓝）、Red = 警示。
% 变量名沿用旧版（Teal / Orange / ...），只改取值，页面内容里的颜色引用不用动。
\definecolor{Ink}{HTML}{0A1A33}
\definecolor{Navy}{HTML}{001E50}
\definecolor{Teal}{HTML}{0054AE}
\definecolor{Blue}{HTML}{0092C2}
\definecolor{Orange}{HTML}{6515AD}
\definecolor{Red}{HTML}{C2334A}
\definecolor{Gold}{HTML}{00377C}
\definecolor{Cyan}{HTML}{00C7FD}
\definecolor{Lav}{HTML}{AF99FF}
\definecolor{Muted}{HTML}{5B6B82}
\definecolor{Line}{HTML}{C9D3E3}
\definecolor{Mist}{HTML}{F5F8FF}
\definecolor{TealBg}{HTML}{E6EFF9}
\definecolor{BlueBg}{HTML}{E0F6FD}
\definecolor{OrangeBg}{HTML}{EFE7F8}
\definecolor{RedBg}{HTML}{FBE9EC}
\definecolor{GoldBg}{HTML}{E4EAF4}
\hypersetup{colorlinks=true,linkcolor=Ink,urlcolor=Blue,citecolor=Blue}

\newcolumntype{Y}{>{\raggedright\arraybackslash}X}
% Pull top-aligned columns closer to the title rule without changing table-only slides.
\newcommand{\minipagetop}{\vspace{-.3cm}}

% ---------- 通用宏 ----------
\newcommand{\memo}[1]{{\color{Muted}\footnotesize #1}}
\newcommand{\pill}[2]{{\setlength{\fboxsep}{1.6pt}\raisebox{.05ex}{\colorbox{#1}{\textcolor{white}{\fontsize{6.5}{7}\selectfont\bfseries #2}}}}}
\newcommand{\tagE}{\pill{Teal}{实验}}
\newcommand{\tagC}{\pill{Blue}{源码}}
\newcommand{\tagO}{\pill{Gold}{官方}}
\newcommand{\tagH}{\pill{Orange}{假设}}
\newcommand{\tagG}{\pill{Red}{未测}}
\newcommand{\labSame}{\pill{Teal}{方向一致}}
\newcommand{\labMag}{\pill{Gold}{量级不同}}
\newcommand{\labBase}{\pill{Orange}{口径不同}}
\newcommand{\labRev}{\pill{Red}{方向相反}}
% 伪代码块：Verbatim 中用 \kw{} \cm{} \id{} 着色（花括号需写成 \{ \}）
\newcommand{\kw}[1]{\textbf{\textcolor{Teal}{#1}}}
\newcommand{\cm}[1]{\textcolor{Muted}{#1}}
\newcommand{\id}[1]{\textcolor{Orange}{#1}}
\newcommand{\hi}[1]{\textbf{\textcolor{Red}{#1}}}
\newtcolorbox{codebox}[2][]{enhanced,colback=Mist,colframe=Line,boxrule=.4pt,arc=2pt,left=4pt,right=3pt,top=2pt,bottom=2pt,
  title={#2},coltitle=white,colbacktitle=Teal,fonttitle=\fontsize{6.8}{8}\selectfont\bfseries,boxsep=1pt,toptitle=1.2pt,bottomtitle=1.2pt,lefttitle=4pt,before skip=0pt,#1}
% 右侧说明框：高度自适应，可连续堆叠
\newtcolorbox{nbox}[3][]{colback=#2!6,colframe=#2,boxrule=.7pt,arc=2.5pt,left=4pt,right=4pt,top=3pt,bottom=3pt,
  before upper={\def\codesize{\relax}\noindent\textbf{\textcolor{#2}{#3}}\par\vspace{1pt}},fontupper=\fontsize{7.7}{9.5}\selectfont\raggedright,before skip=0pt,after skip=4pt,#1}
\newcommand{\hl}[1]{\textbf{\textcolor{Teal}{#1}}}
\newcommand{\warn}[1]{\textbf{\textcolor{Red}{#1}}}
\newcommand{\codesize}{\small}
\newcommand{\code}[1]{\texttt{\codesize #1}}
\newcommand{\tb}[1]{{\bfseries #1}}
% Walk 名字：小字号，长名（如 image_gen_guided3）能放进 2.25 cm 的框
\newcommand{\wname}[1]{{\fontsize{6.6}{8}\selectfont\bfseries #1}}

% 图片来源注记
\newcommand{\figsource}[1]{{\color{Muted}\fontsize{6.5}{7.8}\selectfont #1}}

% ---------- TikZ 图元样式（固定高度、字号，避免重叠） ----------
\tikzset{
  >={Latex[length=1.8mm]},
  box/.style={draw=#1,fill=#1!8,rounded corners=2.5pt,line width=.8pt,align=center,font=\footnotesize,minimum height=.62cm,inner xsep=4pt,inner ysep=2.5pt},
  box/.default=Teal,
  sbox/.style={box=#1,font=\scriptsize,minimum height=.45cm,inner xsep=3pt,inner ysep=1.5pt},
  sbox/.default=Teal,
  plain/.style={align=center,font=\footnotesize,inner sep=1pt},
  tiny/.style={align=center,font=\scriptsize,inner sep=1pt,text=Muted},
  flow/.style={->,draw=Muted,line width=.9pt},
  ctrl/.style={->,draw=Blue,line width=.9pt,densely dashed},
  strm/.style={->,draw=Orange,line width=1.2pt},
  back/.style={->,draw=Teal,line width=1pt,rounded corners=3pt},
  wk/.style={box=#1,minimum height=.62cm,font=\scriptsize,inner xsep=3pt},
  panel/.style={draw=Line,fill=Mist,rounded corners=3pt,line width=.6pt},
}

% 对数轴上的点与连线（第 18 页）
\newcommand{\two}[3]{\pgfmathsetmacro{\xa}{ln(#2/0.5)/ln(2)*1.5}\pgfmathsetmacro{\xb}{ln(#3/0.5)/ln(2)*1.5}\draw[Muted!60,line width=.8pt] (\xa,#1)--(\xb,#1);}
\newcommand{\dotat}[3]{\pgfmathsetmacro{\xx}{ln(#2/0.5)/ln(2)*1.5}\node[#3] at (\xx,#1) {};}
。

\setbeamercolor{normal text}{fg=Ink,bg=white}
\setbeamercolor{structure}{fg=Teal}
\setbeamercolor{alerted text}{fg=Red}
\setbeamertemplate{navigation symbols}{}
\setbeamertemplate{itemize item}{\raisebox{.15ex}{\color{Teal}\scriptsize$\blacksquare$}}
\setbeamertemplate{itemize subitem}{\raisebox{.15ex}{\color{Teal}\tiny$\blacksquare$}}
\setbeamertemplate{enumerate item}{\color{Teal}\bfseries\insertenumlabel.}
\setbeamersize{text margin left=1.0cm,text margin right=1.0cm}
\setbeamerfont{itemize/enumerate body}{size=\small}
\setbeamerfont{itemize/enumerate subbody}{size=\footnotesize}
\setlength{\leftmargini}{1.1em}
\setlength{\parskip}{0pt}
\setlength{\itemsep}{2pt}
\renewcommand{\arraystretch}{1.18}
\setlength{\tabcolsep}{4pt}

\newcommand{\deckpart}{}
% PDF 书签：每个 \setpart 是一级书签，每页标题是二级书签（打开 PDF 时可从侧边栏跳转）。
\hypersetup{bookmarks=true,bookmarksopen=true,bookmarksopenlevel=1,bookmarksnumbered=false,unicode=true}
\newcommand{\setpart}[1]{\renewcommand{\deckpart}{#1}\pdfbookmark[1]{#1}{part.\thepage.\the\numexpr\value{page}\relax}}
\newcommand{\framebookmark}[1]{\pdfbookmark[2]{\thepage\ \ #1}{frame.\thepage}}

\setbeamertemplate{frametitle}{%
  \vspace{-.05cm}%
  \framebookmark{\insertframetitle}%
  \begin{tikzpicture}[remember picture,overlay]
    \fill[Cyan] ([xshift=1.0cm,yshift=-.43cm]current page.north west) rectangle ++(.13cm,-.13cm);
    \node[anchor=west,inner sep=0,font=\fontsize{7}{8}\selectfont\bfseries,text=Muted] at ([xshift=1.25cm,yshift=-.495cm]current page.north west) {\deckpart};
    \node[anchor=north west,inner sep=0,font=\fontsize{14.5}{17}\selectfont\bfseries,text=Ink,text width=14.4cm] at ([xshift=1.0cm,yshift=-.7cm]current page.north west) {\insertframetitle};
    \fill[Cyan] ([xshift=1.0cm,yshift=-1.42cm]current page.north west) rectangle ++(1.0cm,-.045cm);
    \fill[Orange] ([xshift=2.0cm,yshift=-1.42cm]current page.north west) rectangle ++(1.0cm,-.045cm);
    \fill[Teal] ([xshift=3.0cm,yshift=-1.42cm]current page.north west) rectangle ++(1.0cm,-.045cm);
  \end{tikzpicture}%
  \vspace{.85cm}%
}
% 页脚：只保留页码，右下对齐。
\setbeamertemplate{footline}{%
  \begin{tikzpicture}[remember picture,overlay]
    \node[anchor=east,inner sep=0,font=\fontsize{7.5}{9}\selectfont\bfseries,text=Muted] at ([xshift=-1.0cm,yshift=.36cm]current page.south east) {\insertframenumber};
  \end{tikzpicture}%
}

% 页尾要点条：没有标签、不是每页都放；位置贴近页脚，内容区因此更大。
\newcommand{\takeaway}[1]{%
  \begin{tikzpicture}[remember picture,overlay]
    \node[anchor=south west,inner xsep=.38cm,inner ysep=.1cm,minimum height=.74cm,minimum width=14cm,text width=13.3cm,align=left,fill=Mist,font=\fontsize{8.6}{11}\selectfont] (sw) at ([xshift=1.0cm,yshift=.62cm]current page.south west) {#1};
    \fill[Teal] (sw.south west) rectangle ([xshift=.09cm]sw.north west);
  \end{tikzpicture}}

% 封面：白底、大标题、左下蓝色标签、细横线。调用时传入标题、副标题、标签行。
\newcommand{\deckcover}[4]{%
  \begin{frame}[plain,noframenumbering]
    \pdfbookmark[1]{封面}{title.page}%
    \begin{tikzpicture}[remember picture,overlay]
      \fill[Cyan] ([xshift=1.4cm,yshift=-1.35cm]current page.north west) rectangle ++(.2cm,-.2cm);
      \fill[Teal] ([xshift=1.4cm,yshift=-1.55cm]current page.north west) rectangle ++(2.4cm,-.16cm);
      \node[anchor=west,inner sep=0,text width=13.6cm,align=left] at ([xshift=1.4cm,yshift=.35cm]current page.west) {%
        {\fontsize{23}{30}\selectfont\bfseries\textcolor{Navy}{#1}}\\[.5cm]
        {\fontsize{12}{16}\selectfont\textcolor{Teal}{\bfseries #2}}\\[.1cm]
        {\fontsize{10.5}{14}\selectfont\textcolor{Muted}{#3}}};
      \draw[Line,line width=.5pt] ([yshift=1.05cm]current page.south west) -- ([yshift=1.05cm]current page.south east);
      \fill[Teal] ([yshift=.0cm]current page.south west) rectangle ++(2.6cm,1.05cm);
      \node[anchor=center,font=\footnotesize\bfseries,text=white] at ([xshift=1.3cm,yshift=.52cm]current page.south west) {2026-11-23};
      \node[anchor=west,inner sep=0,font=\fontsize{7}{9}\selectfont,text=Muted,align=left] at ([xshift=3.2cm,yshift=.52cm]current page.south west) {#4};
    \end{tikzpicture}
  \end{frame}}


\title{机器人 VLA 实时运行时：输入、推理与当前进展}
\subtitle{框架的异步与流式能力；Pi0.5 与 Qwen3-Omni 两个领域的现状与路线}
\date{2026-11-23}
\hypersetup{pdftitle={机器人 VLA 实时运行时：输入、推理与当前进展}}

\begin{document}

\deckcover{实时多模态推理：框架能力与两个领域的路线\\异步输入、流式输出与全双工}%
  {框架能力 · Pi0.5（机器人 VLA）· Qwen3-Omni（全模态对话）}%
  {M* · vLLM-Omni · SGLang-Omni 的在线与会话能力 · 多频输入 · 测量结果 · 实验计划}%
  {汇报稿\quad|\quad robotics\_vla\_runtime\\版本：M* 2c8ffde8 · vLLM-Omni 5917730+overlay · SGLang-Omni e8c635e · SGLang core 0.5.22}

\setpart{结论}

\begin{frame}[t]{主要结论}
\footnotesize
\renewcommand{\arraystretch}{1.05}
\begin{enumerate}\setlength\itemsep{2.5pt}
\item \tb{输出流式三个框架都有。} M* 的 NDJSON 与分区流式边、vLLM-Omni 的增量事件、SGLang-Omni 的 SSE，都能在请求进行中逐步返回。
\item \tb{输入流式只有会话式入口。} M* 对 Qwen3-Omni 没有输入流（模型未声明，追加被拒绝）。vLLM-Omni 的视频会话与 SGLang-Omni 的会话入口可以持续接收音频和帧，但只在查询或轮次边界时被模型看到。
\item \tb{三个框架都不能把运行中到达的输入并入正在进行的生成。} vLLM-Omni 的新查询抢占旧查询；SGLang-Omni 的打断后重问或截断；M* 没有打断。这是两个领域共同的缺口。
\item \tb{取消与打断各不相同。} vLLM-Omni 的取消会到达引擎的 abort；M* 只丢弃排队中的请求，已启动的一步跑完；SGLang-Omni 的取消未核对。vLLM-Omni 与 SGLang-Omni 有打断，M* 没有。
\item \tb{Pi0.5（机器人）：Level 1 已完成。} 外部运行时在三个框架上测完了多频输入、触发与欠载。要达到目标还需要引擎级改动：同连接接收、中止在途推理、RTC。
\item \tb{Qwen3-Omni：目前只有基准仓库的单请求数据。} 下一步是两条 Level 1 路线（窗口式与会话式），都不改框架；推理中并入新输入属于 Level 2。
\end{enumerate}
\par\takeaway{框架的能力差异在两件事上：输入能否在运行中进入（M* 依赖模型声明，vLLM-Omni 与 SGLang-Omni 依赖会话入口），以及取消能否到达 GPU。}
\end{frame}

\begin{frame}[t]{本汇报的结构}
\scriptsize
\renewcommand{\arraystretch}{1.12}
\setlength{\tabcolsep}{3pt}
\begin{tabularx}{\textwidth}{>{\bfseries}p{1.3cm} >{\bfseries}p{3.4cm} Y}
\toprule
\rowcolor{Mist}
部分 & 主题 & 回答什么 \\
\midrule
一 & 框架能力 & 六项能力的判据；M*、vLLM-Omni、SGLang-Omni 各能做什么、不能做什么；对照；要达到目标各层要改什么 \\
二 & Pi0.5（机器人 VLA） & 输入与输出；多频输入与快照；实时方案（控制与推理解耦、动作块衔接、RTC）；测量；Level 2 的缺口 \\
三 & Qwen3-Omni（全模态对话） & 两种流式；三类输入的时间语义；现有的单请求数据；三层路线与评价指标 \\
四 & 计划与待确定事项 & 实验阶段与状态；需要确定的范围与参数 \\
附 & 版本、来源与局限 & 版本与来源；局限与未验证项 \\
\bottomrule
\end{tabularx}
\par\takeaway{先看一的六项能力，二和三是两个领域的具体问题，四是接下来要做的实验。}
\end{frame}

\setpart{框架能力 · 判据与目标}

\begin{frame}[t]{先定义六种能力，再看框架}
\scriptsize
\renewcommand{\arraystretch}{1.12}
\setlength{\tabcolsep}{3pt}
\begin{tabularx}{\textwidth}{>{\bfseries}p{2.6cm} Y Y}
\toprule
\rowcolor{Mist}
能力 & 指什么 & 框架要具备什么（判据） \\
\midrule
输出流式 & 一次请求进行中，中间结果已经返回（文本 token、音频块、动作块） & 有逐步回传的通道，不必等整段结束 \\
输入流式（运行中追加） & 请求还在运行时，客户端继续送入音频、视频、文本或观测 & 有接收追加输入的协议，并且模型声明了可接收的输入流 \\
推理中并入新输入 & 新输入改变正在进行的计算（上下文、KV、后续输出） & 调度与引擎能把新 token 或帧接到正在进行的序列上 \\
打断（barge-in） & 输出过程中，新输入使当前输出停止或转向 & 有输入事件能触发取消、截断或重新提问 \\
中止在途计算 & 取消后，GPU 上的计算停止 & 取消能进入引擎，并在步与步之间生效 \\
增量处理 & 只计算新增的部分，已有部分复用 & 有前缀复用或增量 prefill \\
\bottomrule
\end{tabularx}
\par\takeaway{输出流式两个领域都已经有，输入流式和"推理中并入新输入"三个框架都还不具备。下面的框架页按这六项判断。}
\par\memo{定义是本汇报的约定，依据是两个领域的目标场景（下一页）。}
\end{frame}

\begin{frame}[t]{两个领域要的能力不同}
\scriptsize
\renewcommand{\arraystretch}{1.12}
\setlength{\tabcolsep}{3pt}
\begin{tabularx}{\textwidth}{>{\bfseries}p{2.6cm} Y Y}
\toprule
\rowcolor{Mist}
能力 & Pi0.5（机器人 VLA） & Qwen3-Omni（全模态对话） \\
\midrule
多路输入独立到达 & 需要：相机 15 / 15 / 7.5 Hz，状态 100 Hz，指令为事件 & 需要：视频 30 帧/秒，音频 20 毫秒一块，文本为事件 \\
窗口式提交（Level 1） & 已完成：外部运行时做快照与触发 & 待做：每次提交一个时间窗口 \\
推理中并入新输入 & 收益有限：Pi0.5 是全注意力，state 写在 prompt 里，任何变化都使前缀失效 & 核心需求：用户说话、画面变化时，正在进行的理解要能吸收 \\
输出流式 & 动作块整块返回；动作分块返回是可选项 & 文本 token 与音频块逐步返回 \\
打断 / 替换 & 新观测到达就替换动作块（重规划） & 用户开口就打断正在播放的回答 \\
中止在途计算 & 需要：限定延迟，替换时停止计算 & 需要：打断时停止生成，不再占用 GPU \\
\bottomrule
\end{tabularx}
\par\takeaway{机器人主要需要"替换"和"限定延迟"，对话主要需要"并入"和"打断"。两个领域对框架的要求不同，不能用同一套评价。}
\par\memo{Pi0.5 的前缀与 state 依据：openpi pi0.py（全部前缀双向）、tokenizer.py:24--28（state 写入文本）；设计文档 §4.2。}
\end{frame}

\setpart{框架能力 · M*}

\begin{frame}[t]{M*：运行中的输入只对声明了输入流的模型开放}
\vspace{-.15cm}
\begin{tikzpicture}[x=1cm,y=1cm,
  s/.style={box=Teal,font=\tiny,align=center,minimum width=1.55cm,minimum height=.62cm},
  ]
  \node[s] (c) at (.8,0) {客户端\\WS 追加帧};
  \node[s] (e) at (2.75,0) {入口\\检查开放与大小};
  \node[s] (d) at (4.7,0) {预处理 worker\\process\_stream\_chunk};
  \node[s] (k) at (6.65,0) {Conductor\\按声明路由};
  \node[s,fill=Gold!14,draw=Gold] (b) at (8.6,0) {消费分区\\StreamBuffer + 分块策略};
  \node[s,fill=Gold!14,draw=Gold] (n) at (10.4,0) {节点输入};
  \draw[flow] (c)--(e); \draw[flow] (e)--(d); \draw[flow] (d)--(k); \draw[flow] (k)--(b); \draw[flow] (b)--(n);
\end{tikzpicture}\\[.05cm]
\scriptsize
\renewcommand{\arraystretch}{1.08}
\setlength{\tabcolsep}{3pt}
\begin{tabularx}{\textwidth}{>{\bfseries}p{4.3cm} Y}
\toprule
\rowcolor{Mist}
情况 & 依据 \\
\midrule
已声明输入流：nemotron\_duplex（16 kHz 音频）、qwen3\_tts（文本）、waypoint（动作） & \code{nemotron\_duplex\_model.py:499--503}；\code{qwen3\_tts\_model.py:725--729}；\code{waypoint\_model.py:277--280} \\
\textcolor{Red}{Qwen3-Omni：未声明输入流，追加被拒绝，请求只能是一次性的} & \code{qwen3\_omni/} 下无 \code{from\_partition=CLIENT}；\code{entrypoint.py:2744--2748}、\code{2765} \\
OpenPI 端点（机器人）：没有追加路径；推理期间不读套接字 & \code{openpi.py} 无 append；\code{rust/server/src/ws/openpi.rs:566, 588} \\
Rust 的 \code{/generate} 路径有追加 & \code{rust/server/src/ws/generate.rs:813--832} \\
\bottomrule
\end{tabularx}
\par\takeaway{M* 的输入流是模型声明的能力，不是框架通用能力。要让 Qwen3-Omni 接收运行中的音频和视频，要在模型侧声明输入流，并实现处理函数。}
\par\memo{来源：external/mstar（2c8ffde8）。接收、路由和缓冲的路径见上图；StreamBuffer 在 \code{mstar/streaming/stream\_buffer.py:71}，数据路由在 \code{conductor.py:2318--2323}。}
\end{frame}

\begin{frame}[t]{M*：输出可以逐步返回；取消只作用于排队中的请求}
\scriptsize
\renewcommand{\arraystretch}{1.1}
\setlength{\tabcolsep}{3pt}
\begin{tabularx}{\textwidth}{>{\bfseries}p{2.2cm} Y}
\toprule
\rowcolor{Mist}
项目 & 实际行为（依据） \\
\midrule
输出流式 & 每个文本 token 或音频块作为一条 NDJSON 返回（\code{EMIT\_TO\_CLIENT}）；Qwen3-Omni 的文本在 \code{qwen3\_omni\_model.py:535--538}，音频块在 \code{738--740}，输出 24 kHz（\code{807--809}） \\
断开即取消 & 客户端断开 $\to$ \code{abort\_submission}（\code{entrypoint.py:2785}）$\to$ 预处理 abort（\code{data\_worker.py:423, 433}）$\to$ \code{conductor.py:1826} \\
\quad 排队中的请求 & 直接删除（\code{conductor.py:1861} 排空参与者后移除） \\
\quad 已在 GPU 上启动的一步 & 照常跑完，结果丢弃（\code{request\_lifecycle.py:206--213}；\code{step\_pipeline.py:1001--1005}） \\
打断（barge-in） & 没有打断消息：\code{/v1/realtime} 把 \code{response.cancel} 当作不支持（\code{serving\_realtime.py:940--945}） \\
\bottomrule
\end{tabularx}
\par\takeaway{M* 的取消作用在请求这一级，不作用在步这一级。这与 Pi0.5 的断开实验一致：断开后，在途推理仍然算完。}
\par\memo{来源：external/mstar（2c8ffde8）。Pi0.5 的实测见设计文档 §2.6.4（M* Rust 前端，15 组）。}
\end{frame}

\begin{frame}[t]{M*：Qwen3-Omni 的图，流式边只在分区之间}
\vspace{-.1cm}
\begin{tikzpicture}[x=1cm,y=1cm,
  p/.style={box=Teal,font=\scriptsize,align=center,minimum width=2.3cm,minimum height=.75cm},
  ]
  \node[box=Muted,font=\scriptsize,minimum width=2.3cm,minimum height=.45cm] (cl) at (1.2,0) {客户端（文本、音频、图像）};
  \node[p] (th) at (1.2,-1.5) {Thinker\\GPU 0--1（TP=2）};
  \node[p] (tk) at (4.6,-1.5) {Talker\\GPU 2};
  \node[p] (cw) at (8.0,-1.5) {Code2Wav\\GPU 2};
  \node[box=Muted,font=\scriptsize,minimum width=2.3cm,minimum height=.45cm] (out) at (8.0,-.0) {音频输出（24 kHz）};
  \draw[flow,densely dashed,Red] (cl)--(th) node[midway,right,font=\tiny,text=Red] {无输入流边};
  \draw[strm] (th)--node[above,font=\tiny]{thinker\_states、thinker\_mask}(tk);
  \draw[strm] (tk)--node[above,font=\tiny]{codec\_tokens}(cw);
  \draw[flow] (cw)--(out);
\end{tikzpicture}
\par\takeaway{M* 的流式在输出侧（分区之间）。输入侧没有流式边。报告 v2 的结论是：一段附件连同它前后的文本是一条 Thinker prefill walk；新附件要开新的 walk，不是并入正在进行的 decode。}
\par\memo{来源：\code{qwen3\_omni\_model.py:761--807}（分区与拓扑；\code{StreamingGraphEdge} 在 \code{mstar/streaming/topology.py:16} 定义）；基准仓库 report v2 08a 章（附件与 prefill walk 的关系，推断，未在本仓库实测）。}
\end{frame}

\setpart{框架能力 · vLLM-Omni}

\begin{frame}[t]{vLLM-Omni：流式视频会话，查询时才被模型看到}
\vspace{.12cm}
\begin{minipage}[t]{5.3cm}\minipagetop
\tb{WS /v1/video/chat/stream}\\[.05cm]
\scriptsize
\begin{tabularx}{\linewidth}{>{\bfseries}p{2.0cm} Y}
\toprule
\rowcolor{Mist}
客户端发送 & 服务端返回 \\
\midrule
\code{session.\allowbreak{}config}（一次） & \code{response.start} \\
\code{video.frame}（JPEG/PNG，可带 frame\_id） & \code{response.\allowbreak{}text.delta} \\
\code{audio.chunk}（PCM16，16 kHz） & \code{response.\allowbreak{}output\_audio.\allowbreak{}delta} \\
\code{video.query}（文本） & \code{response.\allowbreak{}text.done}、\code{session.done} \\
\code{video.done}（结束会话） & \code{video.frames.\allowbreak{}consumed}（帧被用到时） \\
\bottomrule
\end{tabularx}
\end{minipage}\hfill
\begin{minipage}[t]{5.1cm}\minipagetop
\begin{nbox}{Orange}{实际行为（依据）}
\scriptsize
\begin{itemize}\setlength\itemsep{1pt}\setlength\leftmargin{.8em}
\item 只支持 Qwen3-Omni（\code{serving\_\allowbreak video\_\allowbreak stream.py:51--55}）。
\item 服务端缓冲：帧最多 50 个，音频 4 MiB（约 131 秒）。
\item 查询时取缓冲快照；生成中送入的帧\textbf{不进入正在运行的请求}（\code{:465--478}）。
\item 每次查询清空音频（\code{:548--550}）；\code{video.done} 不取消当前查询（\code{:781--787}）。
\end{itemize}
\end{nbox}
\end{minipage}
\par\takeaway{这是窗口式会话：输入持续缓存，理解按查询发生。它不是推理中并入新输入。}
\par\memo{来源：pinned/vllm-omni；文档漏列 ack、frames.consumed、ping，代码中有。}
\end{frame}

\begin{frame}[t]{vLLM-Omni：新查询会抢占旧查询；增量 prefill 有三个条件}
\scriptsize
\renewcommand{\arraystretch}{1.1}
\setlength{\tabcolsep}{3pt}
\begin{tabularx}{\textwidth}{>{\bfseries}p{3.0cm} Y Y}
\toprule
\rowcolor{Mist}
项目 & 实际行为 & 依据 \\
\midrule
查询进行中再发查询 & 不排队、不合并：取消旧查询的任务，旧回答的剩余输出静默丢弃，新查询取当前缓冲快照 & \code{video\_stream\_base.py:437--447, 534--562, 1008--1010} \\
增量 prefill 的条件 1 & 阶段 0（Thinker）开启前缀缓存；pinned 的 \code{qwen3\_omni\_moe.yaml} 默认关闭（latest 默认开） & \code{:84--92}；pinned \code{yaml:32}，latest \code{yaml:40} \\
增量 prefill 的条件 2 & 输出只含文本（modalities 不含 audio）；语音输出时退回旧路径 & \code{:296--298} \\
增量 prefill 的条件 3 & 预热只在没有查询进行时进行；进行中送入的帧，在本次生成结束后才预热 & \code{:465--478} \\
\bottomrule
\end{tabularx}
\par\takeaway{增量 prefill 减少的是查询之间的重复计算，不是生成过程中的输入更新。开启语音输出就失去它。}
\par\memo{来源：pinned/vllm-omni；latest 与 pinned 的差别仅在默认值（pinned 默认关，latest 默认开，CUDA 上）。文档 \code{video\_stream\_api.md:19--21} 写的是 opt-in，与 pinned 一致。}
\end{frame}

\begin{frame}[t]{vLLM-Omni：双工端点中，Qwen3-Omni 是轮次模型}
\scriptsize
\renewcommand{\arraystretch}{1.1}
\setlength{\tabcolsep}{3pt}
\begin{tabularx}{\textwidth}{>{\bfseries}p{2.5cm} Y Y}
\toprule
\rowcolor{Mist}
端点 & 行为 & 依据 \\
\midrule
\code{/v1/\allowbreak realtime}（轮次式） & 输入音频在 commit 前一直缓存；commit 生成一次回答；\code{response.cancel} 调用引擎 abort & \code{realtime/connection.py:355--392, 477--507, 1429--1470} \\
\code{/v1/\allowbreak realtime?\allowbreak duplex=1} & 需要部署配置 \code{session\_mode: duplex}（\code{qwen3\_omni\_duplex.yaml:2--3}）；否则返回错误并以 1008 关闭（文档写的是回退到轮次式，与代码不符） & \code{api\_server.py:\textasciitilde1815--1835}；\code{full\_duplex\_api.md} \\
Qwen3-Omni 插件 & \code{supports\_barge\_in=True}；每个提交的话语是一个新请求；不宣称原生流式 KV；轮次由引擎的 VAD / commit 决定 & \code{plugin.py:370--395} \\
打断过程 & 检测到语音 $\to$ 判定 barge-in $\to$ 结束回答，纪元加一，中止当前请求，旧纪元的输出丢弃，发出 \code{audio.cancelled} & \code{runner.py:2064--2155}；\code{engine\_session.py:1903--1911} \\
\bottomrule
\end{tabularx}
\par\takeaway{vLLM-Omni 的"全双工"在会话层：引擎管理轮次和打断，模型仍按话语出请求。Qwen3-Omni 已接入这个端点。}
\par\memo{来源：pinned/vllm-omni；docs/serving/full\_duplex\_api.md:75--80（支持的模型：MiniCPM-o 4.5、Qwen3-Omni、Nemotron VoiceChat；PersonaPlex 未接入）。}
\end{frame}

\begin{frame}[t]{vLLM-Omni：阶段间的流水与异步输出，只覆盖输出侧}
\vspace{-.1cm}
\begin{tikzpicture}[x=1cm,y=.95cm]
  \node[box=Orange,font=\scriptsize,minimum width=2.2cm,minimum height=.6cm] (a) at (1.1,0) {Thinker\\每个 decode 步 1 个};
  \node[box=Orange,font=\scriptsize,minimum width=2.2cm,minimum height=.6cm] (b) at (4.3,0) {Talker\\codec 每 25 帧（首块 4）};
  \node[box=Orange,font=\scriptsize,minimum width=2.2cm,minimum height=.6cm] (c) at (7.5,0) {Code2Wav\\分块解码};
  \node[box=Muted,font=\scriptsize,minimum width=2.2cm,minimum height=.6cm] (i) at (1.1,-1.5) {输入：一次性 prefill\\（不流式）};
  \draw[strm] (a)--node[above,font=\tiny]{hidden states}(b);
  \draw[strm] (b)--node[above,font=\tiny]{codec}(c);
  \draw[flow,densely dashed] (i)--(a);
  \node[box=Gold,font=\scriptsize,minimum width=3.8cm,minimum height=.6cm] (ao) at (7.0,-1.5) {异步输出：D2H 与打包\\移到后台线程（不同模型）};
\end{tikzpicture}
\par\scriptsize
\begin{itemize}\setlength\itemsep{1pt}
\item \tb{async\_chunk}：阶段间按块交接。文档示例（H800，非本项目测量）：文本生成语音、并发 1 时，首包音频时间从 6.46 秒降到 0.52 秒。
\item 输入不经过这条流水线：第 0 阶段一次性 prefill 完整提示（\code{async\_chunk.md:101--107}）。
\item \tb{取消}：取消会一路中止到 vLLM 调度器（\code{async\_omni.py:381--388}；\code{async\_llm.py:779--781}）；GPU 在下一个引擎步停止（推断，未在 GPU 上测）。
\end{itemize}
\par\takeaway{async\_chunk 让输出更早开始，但输入仍是一次性进入第 0 阶段。}
\par\memo{来源：pinned/vllm-omni（\code{deploy/qwen3\_omni\_moe.yaml:15--22}）；H800 数据来自 async\_chunk.md。}
\end{frame}

\setpart{框架能力 · SGLang-Omni}

\begin{frame}[t]{SGLang-Omni：Qwen3-Omni 的会话入口，按轮次工作}
\scriptsize
\renewcommand{\arraystretch}{1.1}
\setlength{\tabcolsep}{3pt}
\begin{tabularx}{\textwidth}{>{\bfseries}p{2.2cm} Y p{3.7cm}}
\toprule
\rowcolor{Mist}
方面 & 实际行为（依据） & 与目标的关系 \\
\midrule
入口 & \code{/v1/realtime}，需 \code{--enable-realtime}；Qwen3-Omni 会话流水线（\code{cli/serve.py:460--470}，\code{config.py:613--618}） & 可直接作为会话测试入口 \\
输入 & PCM16 音频、摄像头帧、文本；VAD 与 dialog 节点切轮次（\code{session\_node.py:20--27}） & 与 vLLM-Omni 视频会话形式接近 \\
输出 & PCM\_F32 音频 24 kHz；文本与音频的 SSE 流式（\code{openai\_api.py:1048--1060}） & 输出流式已有 \\
打断 & 回答中检测到语音开始 $\to$ INTERRUPT（\code{session\_input.py:262--263}）；未完的话语重问，或截断到客户端已听到的音频（\code{:186--188}）；\code{interaction="turn\_based"}（\code{session\_bridge.py:110}） & 有，但仍按轮次 \\
取消 & 断开 $\to$ \code{coordinator.abort}（\code{coordinator.py:534--540}）$\to$ 引擎 abort（\code{native\_chat\_scheduler.py:482}）；GPU 是否立即停止未核对 & 未核对 \\
并入新输入 & 未见：新输入在话语结束后才成为新一轮 & 未见 \\
\bottomrule
\end{tabularx}
\par\takeaway{SGLang-Omni 与 vLLM-Omni 在输入形式和打断上接近，都按轮次工作；推理中并入新输入两者都未见。}
\par\memo{来源：pinned/sglang-omni（e8c635e）；latest 只抽查了文件存在。本仓库未测会话。}
\end{frame}

\setpart{框架能力 · 对照与缺口}

\begin{frame}[t]{能力对照（一）：输入与输出}
\scriptsize
\renewcommand{\arraystretch}{1.12}
\setlength{\tabcolsep}{3pt}
\begin{tabularx}{\textwidth}{>{\bfseries}p{2.9cm} Y Y Y}
\toprule
\rowcolor{Mist}
能力（Qwen3-Omni 语境） & \textcolor{Teal}{M*} & \textcolor{Orange}{vLLM-Omni} & \textcolor{Blue}{SGLang-Omni} \\
\midrule
输出流式（中间结果返回） & 有：NDJSON，每个 token 与音频块一条 & 有：文本与音频增量事件 & 有：SSE，音频 PCM\_F32 \\
运行中的外部输入（追加） & \textcolor{Red}{无}：Qwen3-Omni 未声明输入流，追加被拒 & 部分：帧与音频缓存，查询时使用 & 部分：会话输入，按轮次切分 \\
输入形式 & 文本；图像（PNG/JPEG/WebP/GIF）；音频与视频（随整段请求提交，采样率未核对） & 音频 PCM16 16 kHz；JPEG/PNG 帧；文本 & 音频 PCM16；摄像头帧；文本 \\
多模态时间语义 & 无 & 部分：帧带 frame\_id，返回 latest\_pts\_ms & 未核对 \\
阶段间流式（输出侧） & 有：\code{StreamingGraphEdge} & 有：async\_chunk & 有（阶段运行时，未逐项核对） \\
\bottomrule
\end{tabularx}
\par\takeaway{三个框架都有输出流式。输入流式只有 vLLM-Omni 与 SGLang-Omni 的会话式入口可用，M* 对 Qwen3-Omni 没有。}
\par\memo{依据见前面三页（M* 第 1--3 页，vLLM-Omni 第 1--4 页，SGLang-Omni 第 1 页）。未核对的格子不作结论。}
\end{frame}

\begin{frame}[t]{能力对照（二）：打断、取消与增量}
\scriptsize
\renewcommand{\arraystretch}{1.12}
\setlength{\tabcolsep}{3pt}
\begin{tabularx}{\textwidth}{>{\bfseries}p{2.9cm} Y Y Y}
\toprule
\rowcolor{Mist}
能力 & \textcolor{Teal}{M*} & \textcolor{Orange}{vLLM-Omni} & \textcolor{Blue}{SGLang-Omni} \\
\midrule
打断（barge-in） & \textcolor{Red}{无}：\code{response.cancel} 被拒绝 & 有：双工端点，Qwen3-Omni 插件 \code{supports\_barge\_in} & 有：会话 INTERRUPT，按轮次 \\
新查询 / 新输入对旧输出的处理 & 未见 & 抢占：取消旧查询，旧输出丢弃 & 截断或重问 \\
中止在途计算 & 部分：排队请求丢弃，已启动的一步跑完 & 有：取消 $\to$ 引擎 abort，步间生效（GPU 停止时机推断） & 部分：断开 $\to$ abort；是否立即停未核对 \\
增量 prefill（只算新增上下文） & \textcolor{Red}{无}（未找到） & 部分：需前缀缓存、只出文本，查询间隙预热 & 未核对（历史保留，\code{thinker\_session.py}） \\
推理中并入新输入（改变正在进行的生成） & \textcolor{Red}{无} & \textcolor{Red}{无} & \textcolor{Red}{无} \\
Qwen3-Omni 可用的入口 & \code{/generate}（窗口式，本仓库未测） & 聊天、视频会话、\code{/v1/realtime}、双工端点 & 聊天、\code{/v1/realtime}（需 \code{--enable-realtime}） \\
\bottomrule
\end{tabularx}
\par\takeaway{三个框架都没有"推理中并入新输入"。打断与取消是 vLLM-Omni 与 SGLang-Omni 已有的，M* 没有。}
\par\memo{依据见前面各页。"推理中并入新输入"这一行是三个框架共同的缺口：未找到任何使正在进行的生成吸收新输入的机制。}
\end{frame}

\begin{frame}[t]{要达到目标，各框架需要改哪一层}
\scriptsize
\renewcommand{\arraystretch}{1.12}
\setlength{\tabcolsep}{3pt}
\begin{tabularx}{\textwidth}{>{\bfseries}p{2.6cm} Y Y Y}
\toprule
\rowcolor{Mist}
目标能力 & M* & vLLM-Omni & SGLang-Omni \\
\midrule
窗口式提交（Qwen3-Omni） & 外部运行时即可（\code{/generate} 逐次提交） & 外部运行时即可（聊天接口） & 外部运行时即可（聊天接口） \\
会话式输入与查询 & 需要模型层改动：声明输入流，实现 \code{process\_stream\_chunk} & 现成（视频会话 API），需要的是音频输出时的增量 prefill & 现成（会话入口），需要验证其 GPU 层面的行为 \\
推理中并入新输入（核心缺口） & 调度与引擎层：把新帧、音频接入正在进行的序列，并处理位置与 KV & 同左，并加上增量编码（音频、视频帧） & 同左 \\
打断到 GPU 停止 & 取消路径要进入 worker 的步级检查 & 已有取消到引擎 abort，需验证打断延迟 & 需验证 abort 是否立即停止批次 \\
Pi0.5：同连接接收新观测（E1） & OpenPI 端点的接收与推理解耦（\code{openpi.py:285}；\code{openpi.rs:566}） & 服务层（\code{connection.\allowbreak{}py:228--278}） & 服务层（\code{ws\_utils.py}） \\
Pi0.5：推理中止（E2） & 去噪步之间加取消点 & 同左（框架内的去噪循环不经过取消路径） & 同左 \\
\bottomrule
\end{tabularx}
\par\takeaway{窗口式不需要改框架；会话式和打断在 vLLM-Omni 与 SGLang-Omni 已经存在；真正的缺口是推理中并入新输入，三个框架都要改引擎。}
\par\memo{改动层的划分：外部运行时（不改框架）→ 协议与服务层 → 调度层 → 模型层 → 引擎层。}
\end{frame}

\setpart{Pi0.5 · 输入与输出}

\begin{frame}[t]{Pi0.5 的输入：我们每次发什么}
\scriptsize
\renewcommand{\arraystretch}{1.05}
\setlength{\tabcolsep}{3pt}
\begin{tabularx}{\textwidth}{>{\bfseries}p{2.6cm} Y p{1.5cm}}
\toprule
\rowcolor{Mist}
输入 & 内容与来源（Level 1 回放） & 频率 \\
\midrule
相机 base\_0\_rgb & 外部相机 1，RGB 图，服务端缩放到 224（来自 exterior\_image\_1\_left） & 15 Hz \\
相机 left\_wrist\_0\_rgb & 腕部相机，RGB 图（来自 wrist\_image\_left，相位错开 22 毫秒） & 15 Hz \\
相机 right\_wrist\_0\_rgb & 外部相机 2，RGB 图（来自 exterior\_image\_2\_left） & 7.5 Hz \\
state & 32 维，归一化：DROID 的 7 维位姿加夹爪，补零到 32；由 15 Hz 插值到 100 Hz & 100 Hz \\
prompt & 语言指令字符串，每个 episode 开始时一次 & 事件 \\
参数 & 去噪步数 10；随机种子 42 & 每次请求 \\
\bottomrule
\end{tabularx}
\par\takeaway{请求内容由框架无关的快照生成，只在最后一步按框架改键名：vLLM-Omni 与 M* 用 \code{observation.images.*}，SGLang core 用 \code{observation.state}。}
\par\memo{来源：src/runtime/dataset.py:22--24；observation.py（StreamSpec、Schedule.snapshot）；exps/tasks/pi05/transport.py:199--228。模型为 pi05\_droid（动作块 15 步）；检查点不带归一化统计，回放的状态是笛卡尔位姿，而 pi05\_droid 训练用的是关节角，所以动作数值没有语义，本项目只测时间。}
\end{frame}

\begin{frame}[t]{Pi0.5 的输出：从输入到一块动作}
\vspace{-.1cm}
\begin{tikzpicture}[x=1cm,y=1cm,
  n/.style={box=Teal,font=\scriptsize,minimum height=.8cm},
  k/.style={box=Orange,font=\scriptsize,minimum height=.8cm},
  ]
  \node[n,minimum width=2.3cm] (img) at (1.0,0) {3 张图像\\{\tiny 缩放到 224，每张 256 token}};
  \node[n,minimum width=2.3cm] (txt) at (1.0,-1.1) {prompt + state\\{\tiny state 离散成 256 个 bin}};
  \node[k,minimum width=2.6cm] (pre) at (3.9,-.55) {前缀 prefill（约 910 token）\\{\tiny 只算一次，写入 KV 缓存}};
  \node[k,minimum width=2.6cm,fill=Red!8,draw=Red] (den) at (7.0,-.55) {动作专家：10 步欧拉\\{\tiny $x \leftarrow x + dt\cdot v$，每步读同一份 KV}};
  \node[n,minimum width=2.0cm,fill=Gold!12,draw=Gold] (out) at (9.9,-.55) {动作块\\{\tiny Level 1：15 $\times$ 32}};
  \draw[flow] (img) -- (pre); \draw[flow] (txt) -- (pre);
  \draw[flow] (pre) -- (den);
  \draw[flow] (den) -- (out);
  \node[anchor=west,font=\tiny,text=Muted] at (0,-2.2) {起点是形状为（batch, horizon, 32）的高斯噪声，t 从 1.0 走到 0；每个请求都重算前缀，跨请求不复用 KV。};
  \node[anchor=west,font=\tiny,text=Muted] at (0,-2.5) {前缀（图像、prompt、state）一变，整段 KV 失效；一次请求只返回一次，没有流式输出（TTFT、ITL 不适用）。};
\end{tikzpicture}
\par\takeaway{一次推理 = 一次前缀计算 + 十次动作专家计算。前缀能否复用，取决于这次请求的图像和 state 是否与上一次相同。}
\par\memo{来源：openpi src/openpi/models/pi0.py:216--279（sample\_actions，默认 10 步）；tokenizer.py:24--28（state 离散化）；pi0.py:105--137（前缀全部双向）；report v2 09a 章；设计文档 §3.1。}
\end{frame}
\begin{frame}[t]{机器人的输入是多频、异步到达的}
\vspace{-.1cm}
\begin{tikzpicture}[x=1cm,y=1cm]
  % 时间轴：0--1.0 秒，8 厘米 = 1 秒
  \def\X{2.3}
  \def\S{8.0}
  % 控制时钟 15 Hz（竖虚线）
  \foreach \i in {0,...,15}{
    \draw[Muted!55,densely dashed,line width=.5pt] ({\X+\i*\S/15},0.45) -- ({\X+\i*\S/15},-2.25);
  }
  % 状态 100 Hz
  \node[anchor=west,font=\scriptsize\bfseries,text=Ink] at (0,0.1) {状态 state};
  \node[anchor=west,font=\tiny,text=Muted] at (0,-0.12) {100 Hz（由 15 Hz 插值）};
  \foreach \i in {0,...,99}{
    \draw[Teal,line width=.5pt] ({\X+\i*\S/100},0.05) -- ({\X+\i*\S/100},0.22);
  }
  % 相机 ext1 15 Hz
  \node[anchor=west,font=\scriptsize\bfseries,text=Ink] at (0,-.5) {外部相机 1};
  \node[anchor=west,font=\tiny,text=Muted] at (0,-.7) {15 Hz，延迟 35 ms};
  \foreach \i in {0,...,14}{
    \fill[Orange!70] ({\X+\i*\S/15-.07},-.62) rectangle ({\X+\i*\S/15+.07},-.42);
  }
  % 腕部相机 15 Hz，相位错开 22 ms
  \node[anchor=west,font=\scriptsize\bfseries,text=Ink] at (0,-.95) {腕部相机};
  \node[anchor=west,font=\tiny,text=Muted] at (0,-1.15) {15 Hz，相位错开 22 ms};
  \foreach \i in {0,...,14}{
    \fill[Orange!70] ({\X+.22+\i*\S/15-.07},-1.07) rectangle ({\X+.22+\i*\S/15+.07},-.87);
  }
  % 外部相机 2 7.5 Hz
  \node[anchor=west,font=\scriptsize\bfseries,text=Ink] at (0,-1.4) {外部相机 2};
  \node[anchor=west,font=\tiny,text=Muted] at (0,-1.6) {7.5 Hz，延迟 45 ms};
  \foreach \i in {0,...,7}{
    \fill[Red!75] ({\X+.4+\i*\S/7.5-.08},-1.52) rectangle ({\X+.4+\i*\S/7.5+.08},-1.32);
  }
  % 指令：事件
  \node[anchor=west,font=\scriptsize\bfseries,text=Ink] at (0,-1.85) {语言指令};
  \node[anchor=west,font=\tiny,text=Muted] at (0,-2.05) {事件，每个 episode 一次};
  \fill[Gold] ({\X-.06},-1.97) rectangle ({\X+.06},-1.77);
  % 时间轴
  \draw[Muted,->] ({\X},-2.45) -- ({\X+\S+.1},-2.45);
  \foreach \tm/\lab in {0/0,0.25/0.25,0.5/0.5,0.75/0.75,1.0/1.0}{
    \draw[Muted] ({\X+\tm*\S},-2.42) -- ({\X+\tm*\S},-2.48);
    \node[font=\tiny,text=Muted] at ({\X+\tm*\S},-2.66) {\lab\,s};
  }
  \node[anchor=west,font=\tiny,text=Muted] at (0,0.62) {灰色虚线：控制时钟，每 66.7 毫秒执行一个动作};
\end{tikzpicture}\\[-.05cm]
\begin{tabularx}{\textwidth}{>{\bfseries}p{3.4cm} Y Y}
\toprule
输入 & 到达的节奏 & 到达延迟（模拟参数，非实测） \\
\midrule
状态（本体感知） & 100 Hz 的高频状态 & 2 毫秒，加 0--1 毫秒抖动 \\
相机与语言指令 & 相机 15 / 15 / 7.5 Hz；指令一次性事件 & 相机 25--45 毫秒，加抖动 \\
\bottomrule
\end{tabularx}
\par\takeaway{各路到达的频率和时刻都不同，而控制器每 66.7 毫秒就要一个动作。所以系统里一定有“最新一条”这种读取方式。}
\end{frame}

\begin{frame}[t]{每次推理读到的是“每路最新一条”，旧数据有多旧}
\vspace{-.1cm}
\begin{minipage}[t]{6.9cm}
\minipagetop
\small\tb{快照时每一路输入的年龄（毫秒，M* 行，阈值 0.5 秒触发）}\\[.1cm]
\begin{tikzpicture}[x=1cm,y=.5cm]
  \def\k{.02}
  \foreach \y/\name/\mid/\hi/\colr in {4/状态/7/12/Teal, 3/相机 1/70/100/Orange, 2/腕部相机/56/90/Orange, 1/相机 2（7.5 Hz）/114/174/Red}{
    \node[anchor=east,font=\scriptsize,inner sep=1pt] at (1.85,\y) {\name};
    \fill[\colr!55] (1.95,\y-.25) rectangle ({1.95+\mid*\k},\y+.25);
    \fill[\colr!18] ({1.95+\mid*\k},\y-.25) rectangle ({1.95+\hi*\k},\y+.25);
    \node[anchor=west,font=\tiny,inner sep=1pt] at ({1.97+\hi*\k},\y) {\mid\ / \hi\ 毫秒};
  }
  \draw[Muted,line width=.4pt] (1.95,.4) -- (1.95,4.6);
\end{tikzpicture}\\[.05cm]
\scriptsize 标注为 p50 / p95。深色：p50；浅色：p50 到 p95。同一时钟采样的对照（lockstep）：最旧输入 p50 为 34 毫秒，p95 为 63 毫秒。
\end{minipage}\hfill
\begin{minipage}[t]{3.4cm}
\minipagetop
\begin{nbox}{Teal}{读法}
\small
最旧的一路是 7.5 Hz 的相机 2：它的采样间隔最长 133 毫秒，延迟 45 毫秒，两者相加决定了快照的新旧。\\[2pt]
动作被执行时，最旧输入的陈旧度 p50 是 406 毫秒（M*，见第 6 部分）。
\end{nbox}
\end{minipage}
\par\takeaway{多频输入的新旧程度由最慢的一路决定，与框架无关：三个框架的快照年龄几乎相同。}
\par\memo{来源：runs/2026-11-20\_\allowbreak{}level1/\allowbreak{}level1\_\allowbreak{}summary.md（快照年龄表，三轮中位数）；设计文档 §2.7 第 4 条。}
\end{frame}

\begin{frame}[t]{三种读取方式：等齐、每路最新一条、近似时间对齐}
\renewcommand{\arraystretch}{1.22}
\scriptsize
\begin{tabularx}{\textwidth}{>{\bfseries}p{2.1cm} Y Y Y}
\toprule
\rowcolor{Mist}
 & 等齐再推理 & 每路最新一条 + 记录年龄 & 近似时间对齐 \\
\midrule
何时推理 & 每一路都到了才推 & 控制需要新动作时推 & 一组输入的时间戳落在窗口内时推 \\
额外等待 & 每次要等最慢一路：7.5 Hz 相机最长 133 毫秒采样间隔，加 45 毫秒延迟 & 无 & 取决于窗口大小 \\
各路时刻 & 基本相同（同一时钟采样，即 lockstep） & 不同，每路记录自己的年龄 & 在时间窗口内 \\
输出频率 & 由最慢的一路决定 & 由控制需要决定 & 不超过最慢的一路（源码推断） \\
主流用在哪 & 数据集的一行记录；离线评测 & LeRobot 的观测槽；本项目的 Level 1 & ROS 2 \code{message\_filters} 的 ApproximateTime，多传感器融合 \\
\bottomrule
\end{tabularx}
\par\takeaway{对实时控制，“每路最新一条 + 年龄”更合适：它不增加等待，年龄可测，超龄的输入可以在运行时拒绝。}
\par\memo{来源：设计文档 §3.3（message\_\allowbreak{}filters 8.0.5，approximate\_\allowbreak{}time.hpp:543–544，无周期触发，无快照）；§4.1 的 Level 1 定义；LeRobot 观测槽 remote\_\allowbreak{}inference/\allowbreak{}engine.py:521–528。}
\end{frame}
\begin{frame}[t]{我们搭的 Runtime：从回放的传感器到机器人}
\vspace{-.15cm}
\begin{tikzpicture}[x=1cm,y=1cm,
  r/.style={box=Gold,font=\scriptsize,minimum width=2.25cm,minimum height=.62cm},
  q/.style={box=Teal,font=\scriptsize,minimum width=2.25cm,minimum height=.62cm},
  ]
  \node[r] (a) at (1.15,0) {多频回放\\{\tiny droid\_100 + 时间表}};
  \node[q] (b) at (3.85,0) {观测快照\\{\tiny 每路 $\le t$ 最新一条}};
  \node[q] (c) at (6.55,0) {触发器\\{\tiny 计数 / 阈值 / 周期}};
  \node[q] (d) at (9.25,0) {后端适配器\\{\tiny 打包 msgpack，WebSocket}};
  \draw[flow] (a) -- (b); \draw[flow] (b) -- (c); \draw[flow] (c) -- (d);
  \node[q,fill=Orange!12,draw=Orange,minimum width=2.25cm] (e) at (9.25,-1.25) {框架（OpenPI 协议）\\{\tiny M* · vLLM-Omni · SGLang core}};
  \node[q] (f) at (6.55,-1.25) {回复解析\\{\tiny 检查 horizon $\times$ 32}};
  \node[q] (g) at (3.85,-1.25) {动作管理\\{\tiny 代数号、对齐、欠载}};
  \node[r] (h) at (1.15,-1.25) {控制时钟 15 Hz\\{\tiny 每 tick 执行一个动作}};
  \draw[flow] (d) -- (e);
  \draw[flow] (e.west) -- (f.east);
  \draw[flow] (f) -- (g);
  \draw[flow] (g) -- (h);
  \node[anchor=west,font=\tiny,text=Muted] at (0,-2.2) {代码：src/runtime/（observation、scheduler、action\_manager、controller），src/backends/（三个适配器），src/metrics/level1.py};
  \node[anchor=west,font=\tiny,text=Muted] at (0,-2.5) {指标：端到端延迟、输入年龄、动作年龄、预定时刻滞后、欠载；单元测试 13 项通过（2026-11-23）};
\end{tikzpicture}
\par\takeaway{Runtime 在框架外面做调度、快照、队列和年龄检查。框架只看到完整的请求，因此可以同样用于三个框架。}
\par\memo{来源：docs/\allowbreak{}design.md §2.7、§4.1；单元测试 tests/\allowbreak{}（假服务器，手算预期，CPU）。}
\end{frame}


\setpart{Pi0.5 · 实时方案}

\begin{frame}[t]{控制循环与推理循环分开运行}
\vspace{-.15cm}
\begin{tikzpicture}[x=1cm,y=.8cm,
  s/.style={box=Gold,font=\scriptsize,minimum width=2.0cm,minimum height=.62cm},
  p/.style={box=Teal,font=\scriptsize,minimum height=.66cm},
  ]
  % 传感器
  \node[s] (c) at (1.0,0) {相机 15 / 7.5 Hz};
  \node[s] (st) at (3.25,0) {本体状态 100 Hz};
  \node[s] (tk) at (5.5,0) {语言指令（事件）};
  % 观测槽
  \node[p,minimum width=7.6cm] (slot) at (3.25,-1.1) {最新观测槽：每路只保留最新一条，记录时间戳，读取时不等待};
  \draw[flow] (c) -- (c|-slot.north);
  \draw[flow] (st) -- (st|-slot.north);
  \draw[flow] (tk) -- (tk|-slot.north);
  % 第三行
  \node[p,minimum width=2.1cm] (trg) at (1.2,-2.4) {触发器\\{\tiny 计数 / 阈值 / 周期}};
  \node[p,minimum width=2.4cm,fill=Teal!14] (req) at (4.0,-2.4) {推理请求\\{\tiny 一次在途，带快照时刻}};
  \node[p,minimum width=2.3cm,fill=Orange!14,draw=Orange] (fw) at (6.7,-2.4) {框架推理\\{\tiny M* / vLLM / SGLang}};
  \node[p,minimum width=2.4cm] (aq) at (9.5,-2.4) {动作队列\\{\tiny 块 + 代数号，旧代数丢弃}};
  \draw[flow] (slot.south) -- (req.north);
  \draw[flow] (trg) -- (req);
  \draw[flow] (req) -- (fw);
  \draw[flow] (fw) -- (aq);
  % 第四行：控制
  \node[p,minimum width=2.4cm,fill=Gold!14,draw=Gold] (ct) at (9.5,-3.9) {控制时钟 15 Hz\\{\tiny 每 66.7 毫秒取一个动作}};
  \node[p,minimum width=2.4cm] (rb) at (1.2,-3.9) {机器人\\{\tiny 欠载时保持最后设定点（假设）}};
  \draw[flow] (aq) -- (ct);
  \draw[flow] (ct) -- (rb);
  \node[font=\tiny,text=Muted,anchor=west] at (3.0,-4.9) {推理慢只让动作变旧或欠载，不会让控制时钟停下};
\end{tikzpicture}
\par\takeaway{推理是异步的：控制时钟不等推理，推理也不等传感器全部到齐。两者只通过“最新观测槽”和“动作队列”交接。}
\par\memo{来源：本项目 src/\allowbreak{}runtime/\allowbreak{}observation.py、scheduler.py、action\_\allowbreak{}manager.py、controller.py；参考做法：LeRobot 观测槽与 ActionQueue（remote\_\allowbreak{}inference/\allowbreak{}engine.py:521--528；policies/\allowbreak{}rtc/\allowbreak{}action\_\allowbreak{}queue.py）。}
\end{frame}

\begin{frame}[t]{动作块怎么衔接：对齐、朴素与欠载}
\vspace{-.1cm}
\begin{tikzpicture}[x=1cm,y=.62cm]
  \def\X{1.75}
  \def\S{4.9}
  \def\dt{0.0667}
  % 时间刻度
  \foreach \tm in {0,0.5,1.0,1.5}{
    \draw[Line,line width=.4pt] ({\X+\tm*\S},-3.7) -- ({\X+\tm*\S},.65);
    \node[font=\tiny,text=Muted] at ({\X+\tm*\S},-3.95) {\tm\,s};
  }
  % 快照与到达
  \draw[Red,densely dashed,line width=.7pt] ({\X+.45*\S},.65) -- ({\X+.45*\S},-3.7);
  \node[font=\tiny,text=Red,anchor=south] at ({\X+.45*\S},.66) {快照 0.45 s};
  \draw[Blue,densely dashed,line width=.7pt] ({\X+.6*\S},.65) -- ({\X+.6*\S},-3.7);
  \node[font=\tiny,text=Blue,anchor=south west] at ({\X+.6*\S+.02},.66) {回复到达 0.6 s};
  % 车道 a：对齐执行
  \node[anchor=west,font=\scriptsize\bfseries,text=Ink] at (0,0) {对齐执行};
  \node[anchor=west,font=\tiny,text=Muted,align=left] at (0,-.42) {LeRobot aligned\\按当前时刻取动作};
  \foreach \i in {0,...,8}{\fill[Teal!60] ({\X+\i*\dt*\S},-.12) rectangle ({\X+(\i+1)*\dt*\S-.02},-.5);}
  \foreach \i in {9,...,14}{\fill[Teal!12] ({\X+\i*\dt*\S},-.12) rectangle ({\X+(\i+1)*\dt*\S-.02},-.5);}
  \foreach \i in {0,...,2}{\fill[Muted!30] ({\X+.45*\S+\i*\dt*\S},-.85) rectangle ({\X+.45*\S+(\i+1)*\dt*\S-.02},-1.23);}
  \foreach \i in {3,...,14}{\fill[Orange!70] ({\X+.45*\S+\i*\dt*\S},-.85) rectangle ({\X+.45*\S+(\i+1)*\dt*\S-.02},-1.23);}
  \node[anchor=west,font=\tiny,text=Muted] at ({\X+1.0*\S+.1},-1.04) {灰：跳过的动作};
  % 车道 b：朴素执行
  \node[anchor=west,font=\scriptsize\bfseries,text=Ink] at (0,-1.9) {朴素执行};
  \node[anchor=west,font=\tiny,text=Muted,align=left] at (0,-2.32) {OpenPI 参考客户端\\从动作 0 开始};
  \foreach \i in {0,...,14}{\fill[Teal!60] ({\X+.6*\S+\i*\dt*\S},-1.95) rectangle ({\X+.6*\S+(\i+1)*\dt*\S-.02},-2.33);}
  \node[anchor=west,font=\tiny,text=Ink] at ({\X+.6*\S+.5},-2.6) {每个动作比计划晚一个回复延迟（150 毫秒）};
  % 车道 c：回复晚到
  \node[anchor=west,font=\scriptsize\bfseries,text=Ink] at (0,-3.2) {回复晚到};
  \node[anchor=west,font=\tiny,text=Muted,align=left] at (0,-3.62) {欠载};
  \foreach \i in {0,...,14}{\fill[Teal!60] ({\X+\i*\dt*\S},-3.2) rectangle ({\X+(\i+1)*\dt*\S-.02},-3.58);}
  \fill[Red!70] ({\X+1.0*\S},-3.2) rectangle ({\X+1.3*\S},-3.58);
  \node[font=\tiny,text=Red] at ({\X+1.15*\S},-3.36) {欠载};
  \foreach \i in {13,14}{\fill[Orange!70] ({\X+1.3*\S+(\i-13)*\dt*\S},-3.2) rectangle ({\X+1.3*\S+(\i-12)*\dt*\S-.02},-3.58);}
\end{tikzpicture}
\par\takeaway{对齐执行避免了朴素执行的固定延迟；但回复一晚，剩下的动作很少，很快又欠载。示意图的时间为说明用，不是测量值。}
\par\memo{来源：设计文档 §2.7（朴素与对齐执行的预定时刻滞后：M* +68 毫秒，SGLang 分桶 +135 毫秒，vLLM-Omni +337 毫秒，对齐下约 0）；action\_\allowbreak{}manager.py 的 locate() 与 execute()。}
\end{frame}

\begin{frame}[t]{主流的几种做法（一）：何时发起推理}
\scriptsize
\renewcommand{\arraystretch}{1.12}
\setlength{\tabcolsep}{3pt}
\begin{tabularx}{\textwidth}{>{\bfseries}p{2.3cm} p{3.3cm} Y}
\toprule
\rowcolor{Mist}
做法 & 项目（版本） & 何时发起推理 \\
\midrule
计数触发 & OpenPI 参考客户端（215abfb） & 每执行固定步数就请求：DROID 满 8 步，LIBERO 每 5 步；与时间和剩余动作无关 \\
回放阈值 & LeRobot v0.7.0（410c095） & 队列剩余的播放时间 $\le$ 取 0.5 秒与（往返时间 + 一个周期）中较大者；pi05 块 50 步，30 fps 时剩 15 步请求 \\
异步 + RTC & openpi05-bsp（dfb7444）；rtc-pi05（HF，MIT） & 按时间：执行到 $t \ge \max(s_{\min}, d)$ 才发起，$d$ 取最近 10 次延迟的最大值；推理在后台线程进行 \\
多传感器对齐 & ROS 2 message\_filters（8.0.5） & 不触发推理；只把传感器消息按时间戳配成一组 \\
\bottomrule
\end{tabularx}
\par\takeaway{计数触发的参考客户端在推理期间阻塞控制循环；回放阈值与异步做法让推理在后台进行，控制循环不等推理。}
\par\memo{来源：agents\_\allowbreak{}results/2026-11-20\_\allowbreak{}openpi-remote-pi05.md 第 2 节；2026-11-20\_\allowbreak{}lerobot-async-rtc.md Q1（execution.py:228--230；factory.py:118）；2026-11-21\_\allowbreak{}level2-research-deep.md A、C 节。}
\end{frame}

\begin{frame}[t]{主流的几种做法（二）：新动作怎么接，能否中止}
\scriptsize
\renewcommand{\arraystretch}{1.12}
\setlength{\tabcolsep}{3pt}
\begin{tabularx}{\textwidth}{>{\bfseries}p{2.3cm} Y Y p{3.0cm}}
\toprule
\rowcolor{Mist}
做法 & 新块怎么接 & 能否中止或修正 & 代价 \\
\midrule
计数触发 & 从动作 0 开始，先进先出；无时间戳，无年龄检查 & 无中止，无 RTC & 每个动作晚一个推理延迟 \\
回放阈值 + 动作队列 & 观测槽只留最新一条；带代数号与游标的队列，对齐替换未执行部分 & 只在客户端丢弃；pi05 默认不开 RTC（需非 chunk 模式） & 回复慢时缓冲不够，动作断档 \\
异步 + RTC & 丢掉已执行的 $d$ 步，新块前缀按已发出的动作做推理期引导 & 服务端仍串行，无取消；openpi05-bsp 的 README 尚无数字 & 去噪每步求导，推理变慢 \\
训练期 RTC & 训练时采样延迟、前缀标为干净；推理时硬替换前缀 & 需要重新训练；EagleVLA 为未来修正模块，机制不同 & 推理开销低，但要训练好的检查点 \\
\bottomrule
\end{tabularx}
\par\takeaway{没有一个开源项目同时具备：服务端在推理中接收新观测、真正中止 GPU 推理，并带 RTC。RTC 都在客户端或模型里实现。}
\par\memo{来源：openpi 的 examples/droid、action\_\allowbreak{}chunk\_\allowbreak{}broker；LeRobot execution.py:204、539--582，configuration\_\allowbreak{}rtc.py:70；rtc-pi05 rtc/\allowbreak{}controller.py:146--183；Kinetix src/\allowbreak{}model.py:267--288。论文数字为作者自报，未复现。}
\end{frame}

\begin{frame}[t]{RTC：让新动作接上已经执行的动作}
\vspace{-.1cm}
\begin{minipage}[t]{5.6cm}
\minipagetop
\scalebox{.92}{\begin{tikzpicture}[x=1cm,y=.55cm]
  \def\X{0.2}
  \def\dd{.7}
  \draw[Line,line width=.4pt] (\X,-3.4) -- (\X+5.1,-3.4);
  \node[font=\tiny,text=Muted,anchor=north east] at (\X+5.1,-3.5) {动作序号};
  % 旧块：已执行（前 d 步）与未执行
  \foreach \i in {0,...,5}{\fill[Muted!45] ({\X+\i*.45},0) rectangle ({\X+\i*.45+.36},.5);}
  \node[font=\tiny,text=Muted,anchor=south] at ({\X+1.1},.55) {旧块已执行的部分};
  % 新块
  \foreach \i in {0,...,9}{\fill[Orange!65] ({\X+\i*.45+1.5},-1.2) rectangle ({\X+\i*.45+1.86},-.7);}
  \node[font=\tiny,text=Orange,anchor=north] at ({\X+3.0},-1.25) {新块（来自新的观测）};
  % 跳变
  \draw[Red,line width=.9pt,<->] ({\X+2.25},.5) -- ({\X+2.25},-.7);
  \node[font=\tiny,text=Red,anchor=west] at ({\X+2.32},-.1) {无 RTC：边界跳变};
  % 冻结的前缀（RTC）
  \fill[Teal!25] ({\X+1.5},-2.0) rectangle ({\X+2.7},-2.5);
  \node[font=\tiny,text=Teal,anchor=west] at ({\X+2.75},-2.25) {前缀 d 步：冻结或引导};
  \node[font=\tiny,text=Muted,anchor=west] at (\X,-2.85) {d = 推理耗时 ÷ 控制周期，向上取整};
\end{tikzpicture}}
\end{minipage}\hfill
\begin{minipage}[t]{5.0cm}
\minipagetop
\scriptsize
\renewcommand{\arraystretch}{1.1}
\begin{tabularx}{\linewidth}{>{\bfseries}p{1.7cm} Y}
\toprule
\rowcolor{Mist}
方式 & 做法与代价 \\
\midrule
无 RTC & 新块直接替换旧块；边界可能跳变 \\
推理期 guided & 前缀误差通过去噪步的 VJP 修正速度场，权重以 $\beta$ 截断（公式见下方）；不需要重训，但去噪步必须可求导 \\
训练期 trained & 训练时随机采样延迟，前缀标为干净动作；推理时硬替换；需要重训 \\
\bottomrule
\end{tabularx}
\end{minipage}
\par\begin{center}\scriptsize 权重 $w_t=\min\big(\frac{t^2+(1-t)^2}{t(1-t)},\,\beta\big)$\end{center}
\par\takeaway{guided 不需要重训，但我们的 M* 与 SGLang core 去噪步都在 CUDA 图里，求导需要重新捕获图；这是引擎级改动。}
\par\memo{来源：Kinetix src/\allowbreak{}model.py:236--258；rtc-pi05 rtc/\allowbreak{}masking.py:122--139；LeRobot configuration\_\allowbreak{}rtc.py:\allowbreak{}41--52，\allowbreak{}modeling\_\allowbreak{}rtc.py:\allowbreak{}46--48；\allowbreak{}设计文档 §0 第 7 条。}
\end{frame}

\setpart{Pi0.5 · 测量结果}

\begin{frame}[t]{同协议下的端到端延迟与单连接可持续周期}
\vspace{-.1cm}
\small\tb{单请求端到端延迟，p50（毫秒），RTX 5880，OpenPI WebSocket，10 步}\\[.08cm]
\begin{tikzpicture}[x=1cm,y=.5cm]
  \def\k{.034}
  \foreach \y/\name/\mid/\colr in {4/M* 默认前端（Rust）/66/Teal, 3/SGLang core（分桶）/97/Blue, 2/SGLang core（默认）/167/Blue, 1/vLLM-Omni/211/Orange}{
    \node[anchor=east,font=\scriptsize,inner sep=1pt] at (2.3,\y) {\name};
    \fill[\colr!70] (2.4,\y-.22) rectangle ({2.4+\mid*\k},\y+.22);
    \node[anchor=west,font=\scriptsize\bfseries,inner sep=1pt] at ({2.42+\mid*\k},\y) {\mid};
  }
\end{tikzpicture}\\[.05cm]
\scriptsize
\renewcommand{\arraystretch}{1.1}
\begin{tabularx}{\textwidth}{>{\bfseries}p{3.4cm} Y Y Y}
\toprule
\rowcolor{Mist}
 & 单连接可持续的最短周期 & 周期 100 毫秒时的表现 & 备注 \\
\midrule
M*（Rust 前端） & 50--75 毫秒 & 每秒 10.0 个请求，无跳过 & 最新版 2c8ffde8 \\
SGLang core（分桶 / 默认） & 约 100 毫秒 / 150--200 毫秒 & 分桶：跳过约 5\% 的触发；默认：排队 & 默认的图按精确长度缓存，只有 4 个名额 \\
vLLM-Omni & 200--300 毫秒 & 每秒 3.1 个请求，跳过 275 个触发 & eager 模式 \\
\bottomrule
\end{tabularx}
\par\takeaway{同一协议下，SGLang 的默认配置不是它的最佳配置。比较时取分桶后的数字，M* 与 vLLM-Omni 的差距在 3 倍左右。}
\par\memo{来源：设计文档 §2.6.1（138 个请求）、§2.6.2（每档 40 个请求）、§2.6.3；runs/2026-11-20\_\allowbreak{}ws-probe/\allowbreak{}；runs/2026-11-21\_\allowbreak{}sglang-buckets-base/\allowbreak{}。}
\end{frame}

\begin{frame}[t]{请求断开不会中止在途推理：新请求等的是剩余时间}
\vspace{-.1cm}
\begin{tikzpicture}[x=1cm,y=.45cm]
  \def\X{2.4}
  \def\S{.075}
  % 时间刻度（毫秒）
  \foreach \tm in {0,50,100,150}{
    \draw[Line,line width=.4pt] ({\X+\tm*\S},-2.5) -- ({\X+\tm*\S},.9);
    \node[font=\tiny,text=Muted] at ({\X+\tm*\S},-2.8) {\tm\,毫秒};
  }
  \node[anchor=west,font=\scriptsize\bfseries,text=Ink] at (0,.5) {请求 A};
  
  \fill[Teal!55] ({\X},.05) rectangle ({\X+66.2*\S},.5);
  \node[font=\tiny,text=white,anchor=center] at ({\X+33*\S},.28) {A：66 毫秒};
  \node[anchor=west,font=\scriptsize\bfseries,text=Ink] at (0,-1.1) {请求 B};
  
  \draw[Red,densely dashed,line width=.7pt] ({\X+26.5*\S},.9) -- ({\X+26.5*\S},-1.6);
  \node[font=\tiny,text=Red,anchor=south west] at ({\X+26.5*\S+.02},-1.45) {B 到达，26.5 毫秒};
  \fill[Orange!55] ({\X+66.2*\S},-1.3) rectangle ({\X+93.2*\S},-1.75);
  \fill[Muted!30] ({\X+26.5*\S},-1.3) rectangle ({\X+66.2*\S},-1.75);
  \node[font=\tiny,text=Ink,anchor=center] at ({\X+46*\S},-1.52) {等 A 的剩余推理};
  \node[font=\tiny,text=Orange,anchor=west] at ({\X+93.2*\S+.05},-1.52) {B 完成：93 毫秒};
  \draw[Blue,densely dashed,line width=.7pt] ({\X+66.2*\S},.9) -- ({\X+66.2*\S},-2.5);
\end{tikzpicture}\\[.05cm]
\scriptsize
\renewcommand{\arraystretch}{1.1}
\begin{tabularx}{\textwidth}{>{\bfseries}p{2.6cm} Y Y Y Y}
\toprule
\rowcolor{Mist}
服务（毫秒，p50） & 对照（单请求） & B，A 关闭连接 & B，A 保持打开 & 预测（不中止） \\
\midrule
M*（Rust 前端） & 66.2 & 93.2 & 93.2 & 105.9 \\
SGLang core & 164.7 & 258.5 & 257.1 & 263.5 \\
vLLM-Omni & 210.4 & 336.4 & 328.2 & 336.7 \\
\bottomrule
\end{tabularx}
\par\takeaway{断开后，服务端仍把 A 算完；如果会中止，B 应接近对照值。测得值与“不中止”的预测相符（SGLang、vLLM 偏差在 3\% 内；M* 低 9--12\%，原因未查明）。}
\par\memo{来源：设计文档 §2.6.4（修正版探测，每框架 15 组，runs/2026-11-20\_\allowbreak{}ws-probe-cancel-v2/\allowbreak{}）。第一版探测曾把 Rust 前端误判为会中止，已修正。}
\end{frame}

\begin{frame}[t]{Level 1：各触发策略下的欠载占比（多频输入，对齐执行）}
\vspace{-.1cm}
\scriptsize
\renewcommand{\arraystretch}{1.15}
\setlength{\tabcolsep}{4pt}
\begin{tabular}{>{\bfseries}p{3.4cm} c c c c}
\toprule
\rowcolor{Mist}
触发策略 & \textcolor{Teal}{M*} & \textcolor{Orange}{vLLM-Omni} & \textcolor{Blue}{SGLang（默认）} & \textcolor{Blue}{SGLang（分桶）} \\
\midrule
周期 100 毫秒 & \colorbox{Mist}{\makebox[1.7cm][c]{0.0\%}} & \colorbox{RedBg}{\makebox[1.7cm][c]{0.5\%}} & \colorbox{Mist}{\makebox[1.7cm][c]{0.0\%}} & \colorbox{Mist}{\makebox[1.7cm][c]{0.0\%}} \\
周期 200 毫秒 & \colorbox{Mist}{\makebox[1.7cm][c]{0.0\%}} & \colorbox{RedBg}{\makebox[1.7cm][c]{3.5\%}} & \colorbox{Mist}{\makebox[1.7cm][c]{0.0\%}} & \colorbox{Mist}{\makebox[1.7cm][c]{0.0\%}} \\
周期 300 毫秒 & \colorbox{Mist}{\makebox[1.7cm][c]{0.0\%}} & \colorbox{RedBg}{\makebox[1.7cm][c]{2.2\%}} & \colorbox{Mist}{\makebox[1.7cm][c]{0.0\%}} & \colorbox{Mist}{\makebox[1.7cm][c]{0.0\%}} \\
每 8 步重规划 & \colorbox{Mist}{\makebox[1.7cm][c]{0.0\%}} & \colorbox{Red!35}{\makebox[1.7cm][c]{21.4\%}} & \colorbox{Mist}{\makebox[1.7cm][c]{0.0\%}} & \colorbox{Mist}{\makebox[1.7cm][c]{0.0\%}} \\
阈值 0.3 秒 & \colorbox{Mist}{\makebox[1.7cm][c]{0.0\%}} & \colorbox{RedBg}{\makebox[1.7cm][c]{4.8\%}} & \colorbox{Mist}{\makebox[1.7cm][c]{0.0\%}} & \colorbox{Mist}{\makebox[1.7cm][c]{0.0\%}} \\
阈值 0.5 秒 & \colorbox{Mist}{\makebox[1.7cm][c]{0.0\%}} & \colorbox{RedBg}{\makebox[1.7cm][c]{2.3\%}} & \colorbox{Mist}{\makebox[1.7cm][c]{0.0\%}} & \colorbox{Mist}{\makebox[1.7cm][c]{0.0\%}} \\
阈值 0.8 秒 & \colorbox{Mist}{\makebox[1.7cm][c]{0.0\%}} & \colorbox{RedBg}{\makebox[1.7cm][c]{5.0\%}} & \colorbox{Mist}{\makebox[1.7cm][c]{0.0\%}} & \colorbox{Mist}{\makebox[1.7cm][c]{0.0\%}} \\
\bottomrule
\end{tabular}
\par\takeaway{M* 与分桶 SGLang 在所有策略下都不欠载。vLLM-Omni 的欠载来自慢回复：块只有 1 秒，一次慢回复就用掉大部分缓冲。}
\par\memo{来源：runs/2026-11-20\_\allowbreak{}level1/\allowbreak{}level1\_\allowbreak{}summary.md（三轮中位数，欠载 tick 占比）；设计文档 §2.7。控制器与输入都是模拟的。}
\end{frame}

\begin{frame}[t]{动作有多新：动作年龄随端到端延迟上升}
\vspace{-.1cm}
\small\tb{周期 200 毫秒，多频输入，对齐执行（毫秒，p50）}\\[.08cm]
\begin{tikzpicture}[x=1cm,y=.4cm]
  \def\k{.017}
  \foreach \y/\name/\e/\a/\colr in {4/M*（Rust）/61/134/Teal, 3/SGLang（分桶）/95/201/Blue, 2/SGLang（默认）/110/266/Blue, 1/vLLM-Omni/227/469/Orange}{
    \node[anchor=east,font=\scriptsize,inner sep=1pt] at (2.3,\y+.12) {\name};
    \fill[\colr!30] (2.4,\y+.05) rectangle ({2.4+\e*\k},\y+.32);
    \node[anchor=west,font=\tiny,inner sep=1pt] at ({2.42+\e*\k},\y+.18) {端到端 \e};
    \fill[\colr!80] (2.4,\y-.32) rectangle ({2.4+\a*\k},\y-.05);
    \node[anchor=west,font=\tiny\bfseries,inner sep=1pt] at ({2.42+\a*\k},\y-.18) {动作年龄 \a};
  }
\end{tikzpicture}\\[.05cm]
\scriptsize
\begin{nbox}{Gold}{朴素执行（OpenPI 参考客户端的做法）}
\scriptsize 每个动作比计划时刻晚约一个回复延迟：M* 晚 68 毫秒，SGLang（分桶）晚 135 毫秒，vLLM-Omni 晚 337 毫秒。以每 8 步重规划为例：朴素执行的 vLLM-Omni 欠载只有 2.3\%，对齐执行时为 21.4\%。朴素执行是用更旧的动作填住缺口换来的。
\end{nbox}
\par\takeaway{动作年龄由触发间隔和端到端延迟共同决定。要降低动作年龄，只能缩短端到端延迟或加快触发；取消旧请求做不到（见断开实验）。}
\par\memo{来源：runs/2026-11-20\_\allowbreak{}level1/\allowbreak{}level1\_\allowbreak{}summary.md；设计文档 §2.7 的表与第 5 条。}
\end{frame}

\setpart{Pi0.5 · Level 2 的缺口与路线}

\begin{frame}[t]{Pi0.5 的 Level 2：七个条目，各自的现状与改动位置}
\scriptsize
\renewcommand{\arraystretch}{1.05}
\setlength{\tabcolsep}{2.5pt}
\begin{tabularx}{\textwidth}{>{\bfseries}p{1.0cm} p{2.2cm} Y p{1.0cm} p{2.0cm}}
\toprule
\rowcolor{Mist}
项 & 目标 & 现状（依据） & 阶段 & 改动位置 \\
\midrule
E1 & 推理中接收新观测，不排队 & 三框架的 OpenPI 端点都是一连接一推理；实测排队时间 = 剩余推理（第 6 部分） & 待做 & 服务层 \\
E2 & 中止在途推理 & 四家断开后仍算完（实测）；vLLM-Omni 的取消路径未用于 Pi0.5 & 待做 & 引擎：去噪步之间加取消点 \\
E3 & guided RTC & 参考实现三个（Kinetix、rtc-pi05、LeRobot），数学一致；框架内没有 & 调研完成 & 引擎：去噪步可求导（CUDA 图需重捕获） \\
E3$'$ & 训练期 RTC & Kinetix、openpi\_rlt、EagleVLA 的做法；需要重新训练 & 调研完成 & 模型 \\
E4 & 按路增量输入 & Pi0.5 全注意力，state 写在 prompt，变化使前缀失效；收益小 & 暂缓 & 模型 \\
E5 & 按路编码器缓存 & 无开源参考；ViT 约占 E2E 的 8\%（插桩数据，上界） & 待测 & 引擎或模型 \\
E6 & 动作块内部流式 & 无参考；OpenPI 也整块返回 & 调研 & 协议 \\
E7 & 跨请求状态 & 不适用：三家不保留跨请求 KV & — & — \\
\bottomrule
\end{tabularx}
\par\takeaway{E1 与 E2 已有同协议实验证据，改动在服务层和引擎层。E3 有参考实现，但框架内没有。E4 与 E7 对机器人的收益有限。}
\par\memo{来源：agents\_results/2026-11-21\_level2-research-deep.md（A--F 节）；设计文档 §4.2、§6、§7。论文数字为作者自报，未复现。}
\end{frame}

\setpart{Qwen3-Omni · 输入与输出}

\begin{frame}[t]{Qwen3-Omni 有两种流式，要分开看}
\vspace{-.1cm}
\begin{minipage}[t]{5.0cm}\minipagetop
\tb{A. 模型内部的输出流（已有）}\\[.05cm]
\scalebox{.95}{\begin{tikzpicture}[x=1cm,y=1cm]
  \node[box=Teal,font=\scriptsize,minimum width=1.45cm,minimum height=.5cm] (t) at (.75,0) {Thinker};
  \node[box=Teal,font=\scriptsize,minimum width=1.45cm,minimum height=.5cm] (k) at (2.6,0) {Talker};
  \node[box=Teal,font=\scriptsize,minimum width=1.45cm,minimum height=.5cm] (c) at (4.45,0) {Code2Wav};
  \draw[strm] (t)--node[above,tiny]{token}(k);
  \draw[strm] (k)--node[above,tiny]{codec 块}(c);
\end{tikzpicture}}\\[.05cm]
\scriptsize 三个分区之间边生产边交接，输出音频分块返回。这是流水线，不是输入的更新。
\end{minipage}\hfill
\begin{minipage}[t]{5.0cm}\minipagetop
\tb{B. 模型外部的输入流（待研究）}\\[.05cm]
\scalebox{.95}{\begin{tikzpicture}[x=1cm,y=1cm]
  \node[box=Muted,font=\scriptsize,minimum width=1.2cm,minimum height=.42cm] (v) at (.6,.3) {视频帧};
  \node[box=Muted,font=\scriptsize,minimum width=1.2cm,minimum height=.42cm] (a) at (.6,-.25) {音频块};
  \node[box=Muted,font=\scriptsize,minimum width=1.2cm,minimum height=.42cm] (x) at (.6,-.8) {文本事件};
  \node[box=Gold,font=\scriptsize,minimum width=1.5cm,minimum height=.5cm] (b) at (2.7,-.25) {缓冲与对齐};
  \node[box=Teal,font=\scriptsize,minimum width=1.4cm,minimum height=.5cm] (th) at (4.55,-.25) {Thinker};
  \draw[flow] (v)--(b); \draw[flow] (a)--(b); \draw[flow] (x)--(b);
  \draw[flow] (b)--(th);
\end{tikzpicture}}\\[.05cm]
\scriptsize 持续到达的多模态输入，在推理期间如何进入同一次计算，是本汇报要回答的问题。
\end{minipage}
\par\takeaway{M* 的 Streaming\allowbreak{}GraphEdge 与 vLLM-Omni 的 async\_chunk 解决的是 A。B 是另一个问题：输入要进入一次正在进行的理解。}
\par\memo{来源：report v2 08a 章（三个分区与边的类型）；设计文档 §2.4（M* 的 StreamingGraphEdge 是分区之间的边）。B 的能力在三个框架中的状态见框架部分。}
\end{frame}

\begin{frame}[t]{Qwen3-Omni 的结构：输入经过编码器进入 Thinker}
\vspace{.05cm}
\begin{tikzpicture}[x=1cm,y=1cm,
  io/.style={box=Muted,font=\scriptsize,minimum width=2.2cm,minimum height=.42cm},
  enc/.style={box=Gold,font=\scriptsize,align=center,minimum width=2.4cm,minimum height=.6cm},
  lm/.style={box=Teal,font=\scriptsize,align=center,minimum height=.62cm},
]
  % 第一行：输入
  \node[io] (tx) at (1.3,0) {文本};
  \node[io] (au) at (4.9,0) {音频};
  \node[io] (im) at (8.5,0) {图像 / 视频};
  % 第二行：编码器（文本直接进入 Thinker）
  \node[enc] (ae) at (4.9,-1.35) {audio\_encoder\\AuT};
  \node[enc] (ve) at (8.5,-1.35) {vision\_encoder\\ViT};
  % 第三行：Thinker
  \node[lm,minimum width=9.0cm,minimum height=.9cm] (th) at (5.2,-2.95) {\textbf{Thinker}：MoE 混合专家 LLM · 30B-A3B · GPU 0--1（TP=2）\\文本、音频、图像 token 在同一个上下文里};
  % 第四行：Talker、Code2Wav、输出
  \node[lm,minimum width=2.9cm] (to) at (1.8,-4.25) {文本输出};
  \node[lm,minimum width=2.9cm] (tk) at (5.0,-4.25) {Talker（语音生成）\\3B-A0.3B · GPU 2};
  \node[lm,minimum width=2.6cm] (cw) at (8.5,-4.25) {Code2Wav（声码器）\\GPU 2};
  \node[io,minimum width=2.2cm] (ao) at (11.4,-4.25) {音频输出};
  % 连线
  \draw[flow] (tx.south) -- (tx.south |- th.north);
  \draw[flow] (au.south) -- (ae.north);
  \draw[flow] (im.south) -- (ve.north);
  \draw[flow] (ae.south) -- (ae.south |- th.north);
  \draw[flow] (ve.south) -- (ve.south |- th.north);
  \draw[flow] (th.south -| to.north) -- (to.north);
  \draw[flow] (th.south -| tk.north) -- (tk.north);
  \draw[strm] (tk.east) -- (cw.west);
  \draw[flow] (cw.east) -- (ao.west);
\end{tikzpicture}
\par\takeaway{输入先经过编码器变成 token，再与文本 token 一起进入同一个 Thinker。所以新的音频或视频不是独立请求，而是要进入 Thinker 的上下文。}
\par\memo{来源：基准仓库 report v2 08a 章（5880 上的 GPU 布局，不是本仓库的设定）。}
\end{frame}

\begin{frame}[t]{三类输入的到达节奏与时间语义}
\scriptsize
\renewcommand{\arraystretch}{1.15}
\begin{tabularx}{\textwidth}{>{\bfseries}p{2.3cm} p{2.6cm} p{2.4cm} Y}
\toprule
\rowcolor{Mist}
输入 & 到达节奏（设定的场景） & 数据形式 & 需要的时间语义 \\
\midrule
视频 & 30 帧/秒 & 单帧图像 & 帧的采集时间（PTS），用于视频的位置编码 \\
音频 & 每 20 毫秒一块（50 块/秒） & PCM 小块 & 每块覆盖的媒体时间区间（起点、终点） \\
文本：用户键入 & 事件，不固定 & 完整的一句 & 事件发生时间 \\
文本：ASR 转写 & 每 300--500 毫秒一段 & 文本片段 & 它对应的音频区间，不是转写完成的时间 \\
文本：系统上下文 & 事件 & 字幕、外部信息 & 生效的时刻 \\
\bottomrule
\end{tabularx}
\par\takeaway{同样是文本，三类的时间语义不同，不能放进同一个队列。ASR 的文本要按它对应的音频区间对齐。}
\par\memo{以上到达节奏是本汇报设定的测试场景，不是测量值。数值来源：设计场景中的 30 FPS 视频、20 毫秒音频块与 ASR 的延迟假设。}
\end{frame}

\begin{frame}[t]{时间窗口：对齐不等于同频率，也不等于每次都等齐}
\vspace{-.1cm}
\begin{tikzpicture}[x=1cm,y=.5cm]
  \def\X{1.9}
  \def\S{8.0}
  \foreach \tm/\lab in {0/0 ms,.5/500 ms,1/1000 ms}{
    \draw[Line,line width=.4pt] ({\X+\tm*\S},-3.0) -- ({\X+\tm*\S},.8);
    \node[font=\tiny,text=Muted] at ({\X+\tm*\S},-3.3) {\lab};
  }
  \node[anchor=west,font=\scriptsize\bfseries,text=Ink] at (0,.3) {视频 30 帧/秒};
  \foreach \i in {0,...,29}{\fill[Orange!70] ({\X+\i*\S/30-.04},.12) rectangle ({\X+\i*\S/30+.04},.42);}
  \node[anchor=west,font=\scriptsize\bfseries,text=Ink] at (0,-.5) {音频 20 毫秒块};
  \foreach \i in {0,...,49}{\fill[Teal!55] ({\X+\i*\S/50},-.7) rectangle ({\X+\i*\S/50+.09},-.4);}
  \node[anchor=west,font=\scriptsize\bfseries,text=Ink] at (0,-1.4) {ASR 文本段};
  \fill[Gold!70] ({\X+.60*\S},-1.6) rectangle ({\X+.88*\S},-1.3);
  \node[font=\tiny,text=Ink] at ({\X+.74*\S},-1.45) {900 ms 到达，对应 600--880 ms};
  \fill[Gold!25] ({\X+.45*\S},-1.6) rectangle ({\X+.6*\S},-1.3);
  \node[font=\tiny,text=Muted] at ({\X+.45*\S},-1.95) {600 ms};
  \draw[Red,line width=.9pt] ({\X+1.0*\S},.9) -- ({\X+1.0*\S},-3.0);
  \node[font=\tiny,text=Red,anchor=south] at ({\X+1.0*\S},.92) {t = 1000 ms：提交一次理解};
  \fill[Red!10] ({\X+0*\S},-2.3) rectangle ({\X+1.0*\S},-2.6);
  \node[font=\tiny,text=Red,anchor=west] at ({\X+.02*\S},-2.45) {窗口 [0, 1000 ms)};
\end{tikzpicture}
\par\takeaway{视频每 33 毫秒一帧、音频每 20 毫秒一块、文本按自己的时间段到达。窗口按媒体时间选取，每路取它覆盖的部分，不要求三路同时到达。}
\par\memo{示意图，时间为说明用，不是测量值。ASR 段的时间区间需要单独记录（音频区间与转写完成时间不同）。}
\end{frame}

\begin{frame}[t]{现有数据：Qwen3-Omni 单请求端到端（基准仓库）}
\scriptsize
\renewcommand{\arraystretch}{1.12}
\setlength{\tabcolsep}{3pt}
\begin{tabularx}{\textwidth}{>{\bfseries}p{3.4cm} Y Y Y}
\toprule
\rowcolor{Mist}
场景（5880，3 卡，公平档） & \textcolor{Teal}{M*} & \textcolor{Orange}{vLLM-Omni} & \textcolor{Blue}{SGLang-Omni} \\
\midrule
文本，1k 提示，256 token，端到端（毫秒，缓存关） & 1,645 & 1,613 & 1,596 \\
语音转文字，端到端（毫秒，缓存关） & 255 & 251 & 278 \\
语音转文字，端到端（毫秒，缓存开） & 252 & 665 & 308 \\
语音合成，每秒音频所需时间（毫秒） & 97 & 153 & 120 \\
\bottomrule
\end{tabularx}
\par\takeaway{这些是单个离线式请求的数据。它们说明三个框架在没有外部输入流时的起点，还不能说明输入更新的能力。}
\par\memo{来源：基准仓库 report v2 08b 章与 01 章，不在本仓库复测。文本与语音转文字为 11 月 11 日批次；缓存开的 M* 文本请求崩溃，未计入。}
\end{frame}

\setpart{Qwen3-Omni · 三层路线}

\begin{frame}[t]{三层路线：每层的边界与三个框架的入口}
\scriptsize
\renewcommand{\arraystretch}{1.1}
\setlength{\tabcolsep}{3pt}
\begin{tabularx}{\textwidth}{>{\bfseries}p{2.4cm} Y Y Y}
\toprule
\rowcolor{Mist}
层 & 定义 & 入口（三框架） & 能回答什么 \\
\midrule
Level 1 外部运行时 & 不改框架。外部程序做多路接收、缓冲、时间对齐，再按窗口或会话提交 & 窗口式：M* \code{/generate}、vLLM-Omni 与 SGLang-Omni 的聊天接口。会话式：vLLM-Omni 视频会话，SGLang-Omni 会话入口；M* 对 Qwen3-Omni 不可用 & 数据管理与时间对齐的代价；每次理解的端到端延迟 \\
Level 2 引擎内输入 & 改框架。运行中到达的音频、帧、文本并入正在进行的生成 & 三个框架都没有。需要调度与引擎层的新机制 & 并入的收益，以及新输入对已有 KV 与位置编码的影响 \\
Level 3 端到端交互 & 改框架与协议。输入与输出并发，支持打断、流式音频输出 & vLLM-Omni 双工端点（Qwen3-Omni 插件，轮次模型）；SGLang-Omni 会话（轮次式打断）；M* 无 & 打断延迟；输出音频的连续性；实时因子 \\
\bottomrule
\end{tabularx}
\par\takeaway{Level 1 能用现有接口测；Level 2 是三个框架都要改的部分；Level 3 在 vLLM-Omni 与 SGLang-Omni 上已有轮次式的版本。}
\par\memo{层次定义与 Pi0.5 的 Level 1 / Level 2 保持一致：Level 1 不改框架，Level 2 改引擎。}
\end{frame}

\begin{frame}[t]{Level 1 的两条实验路线（Qwen3-Omni）}
\scriptsize
\renewcommand{\arraystretch}{1.1}
\setlength{\tabcolsep}{3pt}
\begin{tabularx}{\textwidth}{>{\bfseries}p{2.3cm} Y Y}
\toprule
\rowcolor{Mist}
 & 路线 A：窗口式（每秒提交） & 路线 B：会话式（持续送入） \\
\midrule
输入 & 外部缓冲视频 30 帧/秒、音频 20 毫秒一块、文本事件；每 1 秒取一个媒体时间窗口 & 持续送入帧与音频块；查询或轮次边界时提交理解 \\
提交 & 每秒一个完整请求，三个框架都可以 & vLLM-Omni 用 \code{video.query}；SGLang-Omni 用 VAD 切轮次；M* 不可用 \\
输出 & 每个窗口的文本与音频增量 & 同左；另外记录打断与抢占的情况 \\
用来验证 & 数据管理、时间对齐、引擎的单次延迟 & 会话状态、抢占与打断的行为、增量 prefill 是否生效 \\
\bottomrule
\end{tabularx}
\par\takeaway{路线 A 三个框架都能跑，是对照的基础；路线 B 在 M* 上需要模型侧的改动，因此 M* 只能先做路线 A。}
\par\memo{场景与时间数值为本汇报设定，不是测量。路线 B 的 vLLM-Omni 行为依据见框架部分第 2 页（增量 prefill 只在文本输出时生效）。}
\end{frame}

\begin{frame}[t]{评价指标：两域共用，增加多模态项}
\scriptsize
\renewcommand{\arraystretch}{1.1}
\setlength{\tabcolsep}{3pt}
\begin{tabularx}{\textwidth}{>{\bfseries}p{3.0cm} Y}
\toprule
\rowcolor{Mist}
指标 & 定义与测量方式 \\
\midrule
输入接收延迟 & 媒体产生到框架接收的时间（回放时带采集时间戳） \\
输入排队时间 & 输入在外部缓冲与框架内等待的时间 \\
模态偏差 & 视频、音频、文本的媒体时间偏差（同一窗口内） \\
输入新鲜度 & 模型实际使用的输入距采集的时间 \\
增量 prefill 耗时 & 新输入进入上下文的耗时（路线 B） \\
输入到首个输出 & 最后一个相关输入到第一个输出 token 或音频块 \\
打断延迟 & 用户开口到当前输出停止的时间（Level 3） \\
音频块延迟与实时因子 & 音频输出的分块间隔；每秒音频的处理时间 \\
\bottomrule
\end{tabularx}
\par\takeaway{"输入新鲜度"和"增量 prefill 耗时"是区分外部运行时管理与引擎内输入的指标：前者由外部缓冲决定，后者只有引擎内并入时才有意义。}
\par\memo{指标定义是本汇报的约定；骨架与 Pi0.5 相同，Qwen3-Omni 侧增加模态偏差、输入新鲜度等多模态项。M* 的输出测量与 Pi0.5 相同（EMIT 时间戳）。}
\end{frame}

\setpart{计划与待确定事项}

\begin{frame}[t]{实验计划：每个阶段是否改框架，以及当前状态}
\scriptsize
\renewcommand{\arraystretch}{1.08}
\setlength{\tabcolsep}{3pt}
\begin{tabularx}{\textwidth}{>{\bfseries}p{1.2cm} Y p{1.1cm} p{2.5cm}}
\toprule
\rowcolor{Mist}
阶段 & 内容 & 改框架 & 状态 \\
\midrule
P0 & Pi0.5 基线与协议探测：三框架同协议的端到端、周期、排队、断开、多连接 & 否 & 完成（11-20） \\
P1 & Pi0.5 Level 1：多频回放、七种触发策略、欠载与动作年龄 & 否 & 完成（11-21） \\
P2 & Pi0.5 Level 2 的双连接探测（排队与断开） & 否 & 完成；E1、E2 有证据 \\
P3 & Qwen3-Omni 路线 A：窗口式 Level 1，三个框架（M* 经 \code{/generate}） & 否 & 计划 \\
P4 & Qwen3-Omni 路线 B：会话式 Level 1（vLLM-Omni 视频会话、SGLang-Omni 会话） & 否 & 计划 \\
P5 & 推理中输入的测量：会话中途追加输入，记录是否影响当前输出、抢占与打断延迟 & 否 & 计划 \\
P6 & 引擎改动，依据 P3–P5 的结论选择：M* 的 Qwen3 输入流；vLLM-Omni 的音频增量 prefill；Pi0.5 的 E1、E2 & 是 & 依赖 P3–P5 \\
P7 & 端到端全双工评测：输入与输出并发，打断延迟，实时因子 & 视情况 & 计划 \\
\bottomrule
\end{tabularx}
\par\takeaway{Pi0.5 已完成 Level 1；Qwen3-Omni 从 P3 开始。P6 的引擎改动以 P3–P5 的测量为依据。}
\end{frame}

\begin{frame}[t]{待确定的范围与参数}
\scriptsize
\renewcommand{\arraystretch}{1.12}
\setlength{\tabcolsep}{3pt}
\begin{tabularx}{\textwidth}{>{\bfseries}p{2.6cm} Y}
\toprule
\rowcolor{Mist}
问题 & 影响 \\
\midrule
Qwen3-Omni 的目标输入：30 帧/秒、20 毫秒音频块、每秒一个窗口是否就是目标场景 & 决定 P3--P5 的负载；不同的帧率与块大小结论不同 \\
是否把 SGLang-Omni 纳入正式对比（它有 Qwen3-Omni 会话入口，latest 只抽查过） & 决定 P4 的对照范围 \\
M* 的 Qwen3-Omni 输入流是否在范围内（需要模型侧改动） & 决定 M* 能否进入路线 B \\
数据来源：回放录制，还是真实的音视频采集 & 决定时间戳与抖动是否是测量值 \\
Pi0.5 的 Level 2 优先级：E1（服务层）还是 E2（引擎）先做 & 决定 P6 中 Pi0.5 的第一项 \\
\bottomrule
\end{tabularx}
\par\takeaway{这些问题不阻塞 P3，可以在 P3 的测量中逐步确定。}
\end{frame}

\setpart{附录}

\begin{frame}[t]{版本、来源与口径}
\scriptsize
\renewcommand{\arraystretch}{1.05}
\setlength{\tabcolsep}{3pt}
\begin{tabularx}{\textwidth}{>{\bfseries}p{2.3cm} Y Y}
\toprule
\rowcolor{Mist}
对象 & 版本 & 用在哪里 \\
\midrule
M* & 2c8ffde8（最新版，\code{external/mstar}）；活动报告用的 9c6f1ad 没有 OpenPI 端点 & Pi0.5 的 Level 1 与探测（Rust 前端）；Qwen3-Omni 的代码阅读 \\
vLLM-Omni & 运行 5917730 + overlay（\code{pinned/vllm-omni}，v0.32.0）；最新 2e3ad0a4 只读 & Pi0.5 的 Level 1 与探测；Qwen3-Omni 的代码阅读。latest 的增量 prefill 默认值不同 \\
SGLang-Omni & 运行的检出 e8c635e（\code{pinned/sglang-omni}）；最新 92a4c132 只抽查 & 代码阅读（Qwen3-Omni 会话）；本仓库未运行 \\
SGLang core & pip sglang 0.5.22（\code{envs/sglang-omni}） & Pi0.5 的 Level 1 与探测 \\
LeRobot / OpenPI & v0.7.0（410c095）/ 215abfb；只读 & 参考：观测槽、队列、RTC；协议 \\
Kinetix 等 RTC 参考 & openpi05-bsp dfb7444；rtc-pi05（HF，MIT）；Kinetix 9296f31（MIT） & Level 2 参考 \\
基准仓库报告 & report v2（Qwen3-Omni 的单请求数据） & 第三部分的现有数据 \\
\bottomrule
\end{tabularx}
\par\takeaway{版本不同的数字不能混用。本汇报的 M* 数字来自 2c8ffde8；Qwen3-Omni 的单请求数据来自基准仓库，不是本仓库测的。}
\par\memo{同一台 RTX 5880（GPU 0），主机有其他用户的作业。端到端指客户端发出到收到回复。输入时间表、延迟与抖动是模拟参数，控制器是模拟的 15 Hz。}
\end{frame}

\begin{frame}[t]{局限与未验证的部分}
\scriptsize
\begin{itemize}\setlength\itemsep{1pt}
\item \tb{没有真机，Pi0.5 与 Qwen3-Omni 的输入都不是采集数据。} Pi0.5 的动作数值没有语义，只测时间。
\item \tb{框架代码阅读为主。} 取消到 GPU 的时机（vLLM-Omni、SGLang-Omni）是推断，未在 GPU 上测。
\item \tb{M* 只测了默认 Rust 前端；vLLM-Omni 的长尾中有多少来自主机负载，没有分开。}
\item \tb{文档与代码不一致（vLLM-Omni）：} 文档说 \code{duplex=1} 在非双工部署时回退到轮次式，代码返回错误并以 1008 关闭；文档的协议表漏列了 \code{video.frame.ack}、\code{video.frames.consumed}。
\item \tb{未核对：} SGLang-Omni 的流式与会话的 latest 版本；SGLang 能否在 GPU 上立即停止批次；vLLM-Omni 的音频增量 prefill。
\item \tb{论文数字为作者自报：} A2C2、EagleVLA 的数字未复现。
\end{itemize}
\par\takeaway{要把结论用到真实的音视频与真机上，需要采集数据和 P3--P5 的测量。}
\end{frame}


\end{document}
